\section{Shared results IV: theta pairs, towers and twisted operators}\label{sec:pf:sh-theta}

\providecommand{\shthetaLL}{\mathit{LL}}
\providecommand{\shthetaemb}{\iota}
\providecommand{\shthetaGBig}{\mathsf{GBig}}
\providecommand{\shthetaPol}{\mathcal{P}}
\providecommand{\shthetaloc}{\mathrm{loc}}
\providecommand{\shthetatw}{\mathrm{tw}}
\providecommand{\shthetafd}{\mathcal{F}}
\providecommand{\shthetaRow}{\mathsf{Row}}
\providecommand{\shthetaLinked}{\mathsf{Linked}}
\providecommand{\shthetaGood}{\mathsf{Good}}
\providecommand{\shthetadeg}{\deg}
\providecommand{\shthetaht}{\operatorname{ht}}
\providecommand{\shthetaSB}{\mathcal{S}}

This section proves nineteen results of the formal development \cite{LeanSource2026} which are used by several sections of
Part~II: assembly theorems on linked rows and on the local Galois representations of a datum, the Witt-tower formula for the
constants of seesaw squares, the three square statements on theta integrals, the cuspidality of theta lifts along a grid, and
four statements on doubling integrals, zeta integrals and level vectors. Each result is stated precisely, in the notation of the
sections which use it, and proved: the assemblies and the arguments which they contain (an induction along a Witt tower,
the factorisation of theta integrals through a common quotient, a decomposition of irreducible modules into archimedean and finite
parts, Weil's convergence argument on Siegel sets, the passage from the doubling identity to the Rallis inner product formula, the
computation of a level index) are given in full, and the deeper results of the formal development on which they rest are stated
and cited.

\paragraph{The results and where they are used.}
\begin{itemize}
\item \Cref{sec:pf:sh-theta:rows}: rows give linked rows (\cref{pf:sh-theta:linkedrow}; used in \cref{sec:pf:open0,sec:pf:rows});
the class-sum divisibility row (\cref{pf:sh-theta:classsum}; \cref{pf:open0:geo}); the jump statement, the bounds for canonical hub
data and the admissibility conditions of the first kind at the places above a prime which is not a type prime
(\cref{pf:sh-theta:jump,pf:sh-theta:hub,pf:sh-theta:setpt}; \cref{pf:linkN-adm:sizes,pf:linkN-adm:split} and \cref{sec:pf:linkN-sel}).
\item \Cref{sec:pf:sh-theta:squares}: convergence of theta integrals (\cref{pf:sh-theta:weilconv}), local matching for squares
(\cref{pf:sh-theta:sqlmthm}), the agreement of matched theta integrals and the row property $\sfname{SumOKK}$
(\cref{pf:sh-theta:sqcor}), and the cuspidality of lifts along grids (\cref{pf:sh-theta:cusp}); used in
\cref{pf:rowprops:main,pf:fit:sqtop,pf:fit:sumsize} and \cref{sec:pf:stageU}.
\item \Cref{sec:pf:sh-theta:tower}: the Witt-tower formula $\sfname{Tower}$ and the existence of bottom triples with constants
$\sfname{TriEx}$ (\cref{pf:sh-theta:towerthm,pf:sh-theta:triex}; \cref{pf:fit:sqtop,pf:rowprops:main}).
\item \Cref{sec:pf:sh-theta:doubling}: the unramified doubling integral at inert places (\cref{pf:sh-theta:inert};
\cref{pf:liftinj:euler}), the Rallis inner product formula (\cref{pf:sh-theta:heckeform}; \cref{pf:bundles:thm:chain} and
\cref{pf:fs:stepUp}), and the place models, the finite zeta integral of a theta lift and the formula for the level index times the
norm of the level vector (\cref{pf:sh-theta:placemodels,pf:sh-theta:finzeta,pf:sh-theta:levelindex}; \cref{pf:onto3:adelc}).
\end{itemize}

\paragraph{The main argument: the Witt tower.}
The constant of a seesaw square compares two theta integrals, over the unitary groups of the nodes $(i+1,j+1)$ and $(i,j)$ of a
grid. To compute its dependence on the type $t$, both spaces of the square are replaced by a \emph{theta triple} of stage
$(p,q)$, and hyperbolic planes are removed one at a time: first from the smaller space ($q\to q-1$), then from the larger one
($p\to p-1$), until totally definite spaces are reached. The theta integrals of every stage triple factor through one common
quotient of the space of polynomials (\cref{pf:sh-theta:five}\ref{pf:sh-theta:T2}); this makes the constants of a triple
well defined up to one scalar, and each step multiplies them by a Hecke $L$-value $L(e+1)$ or $L(e)$ and by an explicit factor
which is a ``step multiple'', a quotient of Gauss sums and of divisors of conductors. The induction (\cref{pf:sh-theta:descent})
and the comparison with the bottom triple then give the formula of $\sfname{Tower}$.

\paragraph{Conventions.}
We freely use the notation of the citing sections: grids, Weil data, real matches and square constants
(\cref{sec:pf:rowprops:setting}), the uniformity $\shthetaGBig$ (\cref{pf:fit:uniform}), hub data and admissibility
(\cref{pf:linkN-adm:global}), rows and linked rows (\cref{sec:pf:rows}), doubling pairs (\cref{pf:liftinj:ERdef}) and lift data
(\cref{pf:onto3:adelcdef}). $F=\sfname{ThetaLocalInputs}$ is the bundle of classical results of \cref{sec:families:theta}, and
``Theorem~K'' is the conjunction of $\sfname{ThetaIntegralsAgree}$ and $\sfname{LocalMatching}$ (\cref{def:families:K,thm:theorem-K}).
Results of the formal development which are not proved here are cited by their formal names, with a statement of what is used;
the formal objects are large records, and we describe their components only as far as they are used.

% ---------------------------------------------------------------------------
\subsection{Linked rows and the third-place series}\label{sec:pf:sh-theta:rows}

This subsection proves five results used in the proofs of the open parameter statements and of the Galois link: the
passage from rows to linked rows (used in \cref{sec:pf:open0,sec:pf:rows}), the class-sum divisibility row
(\cref{pf:open0:geo}), and three results on the local Galois representations of a datum at the places above a twisted
prime (\cref{pf:linkN-adm:sizes,pf:linkN-adm:split} and \cref{sec:pf:linkN-sel}): the jump statement, the bounds for
canonical hub data and the seventeen admissibility conditions of the first kind. All five are assemblies of results of the
formal development, which we cite; we state them precisely and give the assembly.

\subsubsection{Linked rows}

We use the rows $\shthetaRow(Q,\mathrm{Lk},H)$, the slots, the parameters $p$ with slot system $p.s$, good slot
$\shthetaGood_p$ and bundle linkage $\mathrm{Lk}_p$, and the linked rows $\shthetaLinked(p,r)$ of
\cref{pf:rows:defrow,pf:rows:deflinked}.

\begin{lemma}[Rows are linked rows]\label{pf:sh-theta:linkedrow}
Let $p$ be a parameter and $H$ a property of classes with raw selections. If $\shthetaRow(\shthetaGood_p,\mathrm{Lk}_p,H)$,
then $\shthetaLinked\bigl(p,(S,\mathrm{Lk})\mapsto\shthetaRow(S,\mathrm{Lk},H)\bigr)$.
\formallabel{GlobalTheory.Fc.RL.LinkedAt.of_row}
\end{lemma}

\begin{proof}
Take the auxiliary slot $R=\mathrm{True}$. The base condition ``built cores satisfy $R$'' holds trivially, eventually in the
scale and for all twisting data (\lref{GlobalTheory.Fc.largeBaseProperty_true}: all exceptional sets empty and all
thresholds $0$), and a row for the slot $\shthetaGood_p$ is a row for the smaller slot $\shthetaGood_p\wedge\mathrm{True}$
(\lref{GlobalTheory.Fc.Row.strengthen}).
\end{proof}

\subsubsection{Class-sum divisibility}

We use the parameters $p(N,S,\Gamma)$ of \cref{sec:pf:open0} (with a norm-choice extension $N$ over the recursion statement
$h_L=\sfname{RecTExists}$, an $S$-extension $S$ and a geometric predicate $\Gamma$), the law extension $N_{\mathrm{law}}$, the
bundles $b$ over $(D,O,c)$ with their comparisons $t$ and parts $b.\mathrm{part}(t)$, and the divisor row
$\mathsf{hDivG}$ (\cref{pf:open0:tab}).

\begin{theorem}[Class-sum divisibility on linked rows]\label{pf:sh-theta:classsum}
Assume the existence of ray class fields ($\sfname{RayField}$). Let $p=p(N,S,\Gamma)$ be a parameter such that every value of
$N$ at a point admits a value of $N_{\mathrm{law}}$, and assume $\shthetaLinked(p,\mathsf{hDivG})$. Then
$\shthetaLinked\bigl(p,(S',\mathrm{Lk})\mapsto\shthetaRow(S',\mathrm{Lk},\mathrm{All}(B))\bigr)$, where $B(b)$ asserts: for every
comparison $t$ of $b$, the part $b.\mathrm{part}(t)$ has $\iota$-point summands, its point has the class-sum pin and satisfies the
compatibility conditions and the additional point condition of the definite side, and the datum $D$ is tied to the size data of
$O$ ($\sfname{RecTied}$).
\formallabel{GlobalTheory.T.class_sum_divisor_all_alt}
\formallabel{GlobalTheory.GeoTies.linked_ties_alt}
\end{theorem}

\begin{proof}
\emph{Exponents.} Since the values of $N$ admit values of the law extension, $N$ has the exponent property of the layers
(\lref{GlobalTheory.ExpTiesOfLayer.hasExp_of_layer_alt}); let $\nu$ be the resulting \emph{extension reading} of the row
position $R_X$ (with the $S$- and $\Gamma$-components of $p$).

\emph{The tie row.} By \lref{GlobalTheory.PlugPt.rowE_part_ties_alt} (from $\sfname{RayField}$ and $h_L$; it contains the
weighted side estimates of the ray field core slot, about $2\,100$ Lean lines), eventually at all bundles linked to the pinned
points of the reading $\nu$, for the slot $\shthetaGood_p\wedge\mathrm{RayPosition}$: the point of the bundle satisfies the
ties $\sfname{TiesB}$, the datum is tied to its size data, and the point satisfies the additional condition $\sfname{More}$. If a
bundle $b$ is linked to a bundle point $\beta$ by $\mathrm{Lk}_p$ (that is, $b$ has a witness over a weighted point $q$), then
every part $b.\mathrm{part}(t)$ is linked to the pinned point of the reading at $q$
(\lref{GlobalTheory.T.exponent_reading_of_hasExp_bundle}). Hence the property ``for all $t$: $\sfname{TiesB}$,
$\sfname{RecTied}$ and $\sfname{More}$ at the part'' holds eventually at all bundles linked by $\mathrm{Lk}_p$ for this slot, and
gives a linked row for it with the auxiliary slot $\mathrm{RayPosition}$, whose base condition holds by
\lref{GlobalTheory.Plug.largeBaseProperty_rayProperty} (from $\sfname{RayField}$).

\emph{Combination.} Combine this linked row with $\shthetaLinked(p,\mathsf{hDivG})$ (\lref{GlobalTheory.Fc.RowL.and}). Let $b$
be a bundle selected by a selection of the combined row; it is built, i.e.\ equal to a bundle with a witness over a weighted
point $q$, and the twist conditions hold at its point (\lref{GlobalTheory.Fc.Ind.pt_realised_and_pb_eq} and the twisting
of the point). For a comparison $t$: the divisor condition of $b$ gives the divisibility statement used by
\lref{GlobalTheory.T.class_sum_divisibility_of_extensionWithLaw} (via \lref{GlobalTheory.Fc.FourBoundsAtBundles.p_div_of_div3g}),
which yields the $\iota$-point summands and the class-sum pin of the part; and $\sfname{TiesB}$, the tie of the datum to its size
data and the realisation of the point give the compatibility conditions
(\lref{GlobalTheory.PlugPt.ties_of_tiesB_alt}). With $\sfname{RecTied}$ and $\sfname{More}$ this is $B(b)$. Conclude by the
implication rule for rows (\lref{GlobalTheory.Fc.Row.imp}), keeping the auxiliary slot.
\end{proof}

\subsubsection{The jump statement}

Let $D$ be a datum and $L$ size data with scale $b=b_L$, positive constants $S_c$ and $\eta=\eta_{\mathrm{sz}}$, and the widths
$w''(L)$ of \cref{pf:linkN-adm:global}. For $1\le j\le m$ and every embedding type $\rho$ the development defines, from the
Harish-Chandra parameters of $D$ at $\rho$, finite sets of rational numbers: the \emph{cloud} $J^{\mathrm{cl}}_{j,\rho}$, the
\emph{low} set $J^{\mathrm{low}}_{j,\rho}$, the \emph{block} set $J^{\mathrm{bl}}$, the set of all jumps $J_{j,\rho}$ and their
twisted and dual versions, together with the centre $c_0=c_0(D)$, the size $\mathrm{sz}_{j,\rho}$ and the bound $h_J(L)$
\cite{LeanSource2026}.

\begin{definition}[Admissible parameters and jump conditions]\label{pf:sh-theta:jumpdef}
The parameters of $L$ are \emph{admissible} ($\sfname{AdmJAlt}(L)$) if: (a) for good families at local indices whose prime is not a
type prime, the local Galois representations are crystalline at the places above the prime, with Hodge--Tate jumps equal to the
level-zero Hodge--Tate weights of the local automorphic data; (b) crystalline is equivalent to de Rham with unramified
Weil--Deligne representation; (c) at the type primes the representations are de Rham with jumps equal to the local Hodge--Tate
weights, and good families have distinct absolutely irreducible constituents; (d) the \emph{dictionary}: for a sign $\epsilon$,
the local Hodge--Tate weights are $\mathrm{sg}\cdot p+c_\tau$ for the parameters $p$ of the datum (with $\mathrm{sg}=\pm1$ according to the
type) and the jump of the rank-one character is $-\mathrm{sg}\cdot\mathrm{sz}-c_0+c_\tau$; (e) the \emph{reading}: the multiplicities
of the jumps of the twisted, dual and derived local lattices are read off from those of the local representations; (f)
$\eta+t+S_c<w''$, $3w''\le b$, $b_T+b'/x\le b$ for all $T$, and $t\le\eta$. The \emph{jump conditions} ($\sfname{JumpConditions}(L)$)
are, for all $1\le j\le m$ and $\rho$:
\begin{enumerate}[label=(J\arabic*),leftmargin=*]
\item\label{pf:sh-theta:J1} for some sign $\epsilon$ the Hodge--Tate jumps of the local modules are given by the jump sets;
\item\label{pf:sh-theta:J2} $\mathrm{sz}\le q-c_0\le\mathrm{sz}+S_c$ on the cloud, $(t+3)/2\le|q-c_0|\le\eta$ on the low set, $|q-c_0|\le c_0$ on
the block set, and the cloud, low and full sets have $k-1$, $m$ and $n_V$ elements;
\item\label{pf:sh-theta:J3} low elements lie in $[c_0-\eta,-2]\cup[t+1,c_0+\eta]$, $-1$ and $t$ are not jumps, cloud elements exceed
$t+1$, and $0$ is neither a cloud nor a low element;
\item\label{pf:sh-theta:J4} the low set has $m-j$ negative and $j$ positive elements, and the twisted sets have the corresponding
sign counts $j$, $m-j$, $m-j$ and $k-1+j$;
\item\label{pf:sh-theta:J5} $|q|\le b+S_c+\eta+1+c_0\le h_J$ for all jumps $q$, and $b\le h_J\le2b$;
\item\label{pf:sh-theta:J6} the dual cloud elements have absolute value in $[\mathrm{sz}-\frac{t+1}2,\mathrm{sz}+\frac{t+1}2+S_c]$, the
dual low elements lie in $[-\eta-\frac{t+1}2,\eta-\frac{t+1}2]$, $0$ is not a dual jump, and the negative dual jumps number $m-j$,
resp.\ $k-1+j$;
\item\label{pf:sh-theta:J7} if $\rho$ is a thick place, the dual cloud elements have absolute value in
$[b-\frac{t+1}2,b-\frac{t+1}2+\eta+t+2+S_c]$, and dual low and block elements have absolute value at most $\eta+t+2$.
\end{enumerate}
$\sfname{JumpClosed}(D)$ asserts $\sfname{AdmJAlt}(L)\Rightarrow\sfname{JumpConditions}(L)$ for all size data $L$ of $D$.
\formal{Formalism.MainNotions.AdmissibleParameterJ.AdmJAlt}\formal{Formalism.MainNotions.AdmissibleParameterJ.JumpConditions}
\formal{Formalism.MainNotions.AdmissibleParameterJ.JumpClosed}
\end{definition}

\begin{theorem}[The jump statement]\label{pf:sh-theta:jump}
$\sfname{JumpClosed}(D)$ holds for every datum $D$.
\formallabel{GlobalTheory.JumpStatement.holds_all_alt}
\formallabel{GlobalTheory.JumpSets.stmtJ_holds_alt}
\end{theorem}

\begin{proof}
Let $L$ have admissible parameters. By (f), the \emph{real place facts} hold for all $1\le j\le m$ and $\rho$
(\lref{GlobalTheory.JumpSets.jFacts_of}): the cloud parameters $v_a$ are non-negative, at most $S_c$ and strictly
decreasing, the block parameters $b_i$ are strictly increasing, $\mathrm{sz}_{j,\rho}-b_i\ge(t+3)/2$ for $i<j$ and $\le-(t+3)/2$ for
$i\ge j$, $|\mathrm{sz}_{j,\rho}-b_i|\le\eta$, $\eta+t+S_c<\mathrm{sz}_{j,\rho}\le b+\eta+1$, $b\le\mathrm{sz}_{j,\rho}$ at thick places,
$\eta+t+S_c\le b$ and $t\le\eta$. (For the inequality $\eta+t+S_c\le b$: by (f), $\eta+t+S_c<w''$ and $3w''\le b$, and $w''>0$.)
Then (J1) follows from (a)--(f) by the dictionary and the reading
(\lref{GlobalTheory.JumpSets.jumpLemma_a_of_admissible_alt_alt}), and (J2)--(J7) follow from the real place facts by counting
and comparing the explicit elements of the jump sets: (J2), (J3), (J4), (J5), (J6) and (J7) are, in this order,
\lref{GlobalTheory.JumpSets.jumpSetBound_of_realPlaceFacts}, \lref{GlobalTheory.JumpSets.jumpSetCloud_of_realPlaceFacts}, \lref{GlobalTheory.JumpSets.jumpSetDistinct_of_realPlaceFacts}, \lref{GlobalTheory.JumpSets.jumpSetEta_of_realPlaceFacts}, \lref{GlobalTheory.JumpSets.jumpSetFinal_of_realPlaceFacts}, and \lref{GlobalTheory.JumpSets.jumpSetGap_of_realPlaceFacts}
(the proof of (J5) uses in addition $\eta+t+S_c\le b$, $t\ge1$, $S_c>0$ and $\eta>0$).
\end{proof}

\subsubsection{Canonical hub data}

We use the hub data, their adaptation $\sfname{DataAD}$ to $L$ and the row constants, the setup conditions and the admissibility
conditions of \cref{pf:linkN-adm:def-hub,pf:linkN-adm:def-setup,pf:linkN-adm:def-adm}, and write $\mathcal G=\{g_0,g_1\}$ for the
type primes. For a local index $\iota$ let $\mathrm{RestPt}(\iota,\pi)$ denote the \emph{rest point conditions} (good family,
$\rho_{\iota,\pi}$ is the representation of $\pi$, semisimplicity, continuity, the tame shape, the cyclotomic inertia image and the
non-genericity conditions of the formal definition).

\begin{theorem}[Bounds for canonical hub data]\label{pf:sh-theta:hub}
Assume Fontaine's $B_{\mathrm{cris}}$, local finiteness and Chebotarev's theorem. Let $D$ be a datum with size data $L$, local
Hodge--Tate data $\mathrm{HW}$, Fontaine data $F$, row constants $R$, a norm choice $C$, hub data $X$ which are canonical for
$(L,\mathrm{HW},R)$, $c\ge120$ and $C_T\in\N$ with $C_T\le b_L$, and assume a parameter choice, $\Lambda\ge1$, admissible parameters,
polynomially bounded parameters, and the following properties of $D$ (fields of the norm-choice facts): admissible sums,
provision $41$, the conclusions $42$, $G_5$ and $Z$, residual irreducibility in rank $r$, the $\OO$-splitting, the coefficient
duties, the norm-choice tie, good triples, the same-twist condition, the block character parameters, unramifiedness of $\mu$ and
the depth-zero condition. Then:
\begin{enumerate}[label=(H\arabic*),leftmargin=*]
\item the families of $X$ are good, and the rest point conditions hold at all $\pi\in Z^{\mathrm{good}}_\iota$ for $\ell_\iota\in\{q,d\}$
and for $\ell_\iota\notin\mathcal G$;
\item the Fontaine degrees of the quotients at the distinguished primes, and the adaptation $\sfname{DataAD}$ of $X$;
\item at every twisted $\iota$ with $\ell_\iota\notin\mathcal G$ in the range of $X.\mathrm{exc}$, the block bound
$c_B(\iota)\le c\,tnn_V\Lambda$, and, if $\ell_\iota$ is split, the column bound $c_0(\iota,u)\le c\,nn_V\Lambda$ for every column $u$;
\item Hodge--Tate bounds $(M_0,h_0)$ for the wall condition $\sfname{JointCond\_H0Wall}$, and for every $\iota$ in the range of
$X.\mathrm{exc}$ and every $\pi$ Hodge--Tate bounds for the tame bound condition with $X.c_B(\iota)$;
\item at every twisted split $\iota$ in the range with $\ell_\iota\notin\mathcal G$: the declaration of $X$ is valid for every
choice of embedding types given by (J1); the places of $E$ above $\ell_\iota$ are unramified over $\Q$; the ambient type
conditions at the conjugate places; for every $\pi\in Z^{\mathrm{good}}_\iota$ the column point data (unramified triviality of
$H^1_f$, the twisted dual, $T_2$, Wach and filtration conditions, finiteness of $H^1_f$ up to torsion of finite rank, the condition
$\mathrm{adm}_{51}$, and the decomposition invariants condition with Hodge--Tate bounds); the local fields $K_w$ lie in the maximal
unramified extension; and the localised duals are crystalline for good coefficient extensions;
\item the signature choice condition for $(X.\rho,C)$;
\item at twisted $\iota$ with $\ell_\iota=g_k$ a type prime: good families, $\rho$ the representation of $\pi$ with semisimple
twist for $\pi\in Z_\iota$, and $s,S,\mathrm{sg},U$ with $\eta_X+t+1<s$, $C_T((s+S)(n_V+1)^2+(n_V+1)^3)<g_k$, the sign condition and
the host jump condition;
\item $\eta_X\le2\eta_{\mathrm{sz}}$.
\end{enumerate}
\formallabel{GlobalTheory.T.split_bound_of_params_alt}
\formallabel{GlobalTheory.HubX.Assembly.main_alt_alt}
\end{theorem}

\begin{proof}
The jump conditions hold by \cref{pf:sh-theta:jump} (applied through \lref{GlobalTheory.T.jump_closed_forall_d7}).
The families of $X$ are good (a field of canonical hub data), and the rest point conditions at $q,d$ and off $\mathcal G$ are
\lref{GlobalTheory.restPt_atQD_alt} and \lref{GlobalTheory.restPt_of_ptFam_aux41_oSplit} (from $B_{\mathrm{cris}}$, provision
41 and the $\OO$-splitting). The Fontaine degrees are \lref{GlobalTheory.T.fontaine_quotient_degrees_of_ptFam}; adaptation is a
field of canonical hub data. (H3) is \lref{GlobalTheory.HubX.blockAndTorsionExponentBounds_canonicalHubData} (the canonical
exponents are explicit, and $c\ge120$). (H4): \lref{GlobalTheory.HubX.jointCond_H0Wall_of_canon_alt} and the range bound
\lref{GlobalTheory.T.joint_conditions_range_bound_alt} (from $B_{\mathrm{cris}}$, local finiteness, the jump conditions for all
size data with admissible parameters, polynomially bounded parameters, provision 41, the $\OO$-splitting and the exponent bound
$c_B\ge c_{B,\mathrm{num}}$ of canonical data). (H5): the split row prime and the split declaration
(\lref{GlobalTheory.HubX.SplitDecl.splitRowPrime_of} and \lref{GlobalTheory.HubX.SplitDecl.main}), the ambient type at conjugate places
(\lref{GlobalTheory.T.ambientType_conjugatePlace_at_splitPrime}), unramified crystallinity at the range prime
(\lref{GlobalTheory.HubX.unrCrysAtRangePrime_alt}), the column point data
(\lref{GlobalTheory.T.unramifiedTrivialH1f_of_decl}, \lref{GlobalTheory.T.ambT2_of_deRham_family}, \lref{GlobalTheory.T.admissibleAtGoodPoints_of_coefRing_taken_at_live}, and \lref{GlobalTheory.HubX.jointCondH0all_canon_alt}
and the range bound). (H6) is \lref{GlobalTheory.JointSix.main} (from Chebotarev's theorem, the norm-choice tie, good triples,
the same-twist condition, the good point row \lref{GlobalTheory.Proof.LocalGlobalAssembly.main} and residual absolute
irreducibility). (H7) and (H8) are \lref{UpperTheory.GalS.hubX_typePrime_conditions_alt} (from the parameter choice, the jump
conditions, polynomially bounded parameters, conclusion $G_5$, the absence of a local-global condition
\lref{UpperTheory.GalS.noLgc_alt}, the depth-zero condition and $C_T\le b_L$).
\end{proof}

\subsubsection{Admissibility conditions of the first kind}

\begin{theorem}[Admissibility of the first kind off the type primes]\label{pf:sh-theta:setpt}
Assume $\sfname{Hypotheses}_2$. Let $D$ be a datum with size data $L$, a parameter choice, a norm choice $\mathrm{nch}$, local
Hodge--Tate, Fontaine and comparison data, hub data $X$ and $C_2\in\R$ satisfying the setup conditions
(\cref{pf:linkN-adm:def-setup}). Then for every local index $\iota$ whose prime $\ell_\iota$ is not a type prime and every
$\pi\in Z^{\mathrm{good}}_\iota$, the seventeen admissibility conditions of the first kind of \cref{pf:linkN-adm:def-adm} hold at
$(\iota,X.r_\iota,X.\rho_\iota,\pi)$.
\formallabel{GlobalTheory.T.set_pt_of_setup_not_twisted}
\formallabel{GlobalTheory.SetPtOff.main_alt_alt}
\end{theorem}

\begin{proof}
By \cref{pf:sh-theta:hub}, applied with the row constants $\mathcal K(C_2)$, $c=120$ and $C_T=0$ (its hypotheses are fields of the
setup and of $\sfname{Hypotheses}_2$, and $\Lambda\ge1$), the rest point conditions hold at $(\iota,\pi)$. The class of $\pi$ is
represented in $Z_\iota$, and the jump conditions hold (\cref{pf:sh-theta:jump}). The rest point conditions give the semisimple
decomposition $r_{\iota,\pi}\otimes K_\iota\cong(\rho_{\iota,\pi}\oplus\bigoplus_c\chi_{\iota,c})\otimes K_\iota$
(\lref{Formalism.MainNotions.N.semisimple_decomposition_of}). The seventeen conditions are then: (G1)
\lref{Formalism.MainNotions.N.galois_condition_one_of_alt}; (G2) \lref{GlobalTheory.GTwoTheta.main_alt_alt} (from
$B_{\mathrm{cris}}$ and the jump conditions); (G3) \lref{Formalism.MainNotions.N.galois_condition_three_of_alt_alt} (with the
Burnside statement \lref{GlobalTheory.SetPtOff.litBurnside_all}); unramifiedness outside $\ell$, $S_\iota$, $Q_\iota$
(\lref{Formalism.MainNotions.N.admG4_of_goodR}); the tame shape (\lref{Formalism.MainNotions.N.galois_condition_five_of});
the tame generator conditions (\lref{GlobalTheory.TameDesign.main}); the eigenvector criterion
(\lref{Formalism.MainNotions.N.galois_condition_six_of}); local Tate duality for $H^2$ and the Bloch--Kato conditions
(\lref{GlobalTheory.T.blochKato_localCondition_all_localizations}); continuity and semisimplicity
(\lref{Formalism.MainNotions.N.continuous_family_of} and \lref{Formalism.MainNotions.N.semisimple_family_of}); the decomposition (above);
the cyclotomic inertia image (a rest point condition); compatibility of local coefficients
(\lref{GlobalTheory.T.completion_comparison_admissible_every_place_over}); the inclusion
$K_\lambda\subseteq\Q_\ell(\mu_{\lcm(\ell,N)})$ (a rest point condition); Fontaine's properties
(\lref{GlobalTheory.FontaineAdm.fontaineProperties_all_places}); the integral $H^2$-duality
(\lref{GlobalTheory.T.integral_tate_H2_top}); and the bound by $v_\ell(h_E)$ for unramified homomorphisms
(\lref{GlobalTheory.T.unramifiedClassFieldBound_of_artin}, from Artin reciprocity).
\end{proof}

\begin{remark}\label{pf:sh-theta:setptrem}
The formal name of \cref{pf:sh-theta:setpt} speaks of places which are ``not twisted'', but its hypothesis is that the prime
$\ell_\iota$ is not one of the two type primes; no twisting condition is imposed on $\iota$.
\end{remark}

% ---------------------------------------------------------------------------
\subsection{Theta integrals of squares: convergence, local matching and cuspidality}\label{sec:pf:sh-theta:squares}

This subsection proves the three \emph{square statements} used in \cref{pf:rowprops:main} and in the proof of
\cref{pf:fit:sqtop} (convergence of theta integrals, the local matching of the two Fock spaces of a square, and the
agreement of matched theta integrals), and the cuspidality of theta lifts along the arrows of a grid. We use the notation
of \cref{def:pf:rowprops:weil}: for hermitian spaces $M,X$ over the CM field $E$ of a datum, a sign $\eps$, a standard
additive character $\psi$ and archimedean data $\theta$, a \emph{Weil representation} $\omega$ of the pair $(\eps M,X)$
with splitting characters $(\chi,\chi')$ is a representation of $\Uni(M)(\A)\times\Uni(X)(\A)$ on the space of
\emph{theta functions} on the adelic symplectic space $(\eps M)\otimes X$, i.e.\ of continuous functions $\Phi$ which are
quasi-periodic under the rational points ($\Phi(\gamma+w)=\psi(-\langle\gamma,w\rangle)\Phi(w)$) and invariant under
Heisenberg translations by an open subgroup of the finite part; the \emph{theta kernel} is
$\theta(g,h;\Phi)=(\omega(g,h)\Phi)(0)$. The height $\shthetaht(h)\ge1$ of $h\in\Uni(X)(\A)$ is the product of the
archimedean and the finite local heights.

\subsubsection{Convergence}

The statement $\sfname{SiegelSetUnitary}$ is the reduction theory of unitary groups in the following form
\cite{Borel1969,MoeglinWaldspurger1995}: for a CM extension $E/F$, a non-degenerate hermitian matrix $J$ of size $n$ and a
\emph{Witt frame} $(X_f,Y_f)$ of rank $r$ (bases of complementary maximal isotropic subspaces with ${}^t\bar X_fJY_f=1$ and
anisotropic orthogonal complement), there are a compact subgroup $K$, a compact subset $\Omega$ of the minimal parabolic
subgroup and $c>0$ such that, with the \emph{cone} $\mathrm{cone}(c)=\{t\in\R_{>0}^r:\ ct_{i+1}\le t_i,\ c\le t_r\}$ and the
torus elements $a_t$ (acting by $t_i$ on $X_f$, $t_i^{-1}$ on $Y_f$ and trivially on the complement at the real places), the
\emph{Siegel set} $\mathfrak S=\{pa_tk:\ p\in\Omega,\ t\in\mathrm{cone}(c),\ k\in K\}$ is closed, every element of the adelic group
is $\gamma s$ with $\gamma$ rational and $s\in\mathfrak S$, the conjugates $a_t^{-1}pa_t$ ($p\in\Omega$, $t\in\mathrm{cone}(c)$)
lie in a fixed compact set $C_0$, and for every Haar measure $\mu$ there is $C_1$ with
\[
\int_{\mathfrak S}f\,d\mu\le C_1\int_{\mathrm{cone}(c)}H(t)\,\bigl(\delta(t)\textstyle\prod_it_i\bigr)^{-1}\,dt,\qquad
\delta(t)=\prod_{i=1}^rt_i^{2[F:\Q](n-2i+1)},
\]
for all measurable $f,H\ge0$ with $f(pa_tk)\le H(t)$; it also contains a comparison of heights and the rapid decay of cusp
forms on $\mathfrak S$, which are not used here.

\begin{theorem}[Convergence of theta integrals]\label{pf:sh-theta:weilconv}
Assume $\sfname{SiegelSetUnitary}$. Let $\omega$ be a Weil representation of $(\eps M,X)$, $p\in\N$ such that the form of
$M$ has signature $(\rank M-p,p)$ at some real place and $\rank X<2(\rank M-p)$, $\nu$ a Haar measure on $\Uni(M)(\A)$
with a fundamental domain $\shthetafd$ for $\Uni(M)(F)$, and $\Phi$ a theta function. Then there are $N\in\N$ and $C\in\R$
such that for every $h\in\Uni(X)(\A)$ the function $g\mapsto\theta(g,h;\Phi)$ is integrable on $\shthetafd$ and
\[
\int_{\shthetafd}|\theta(g,h;\Phi)|\,d\nu(g)\le C\,\shthetaht(h)^N .
\]
That is, $\sfname{SquareSt.WeilConvergence}$ holds; in particular $\sfname{SiegelSetUnitary}$ and $F$ imply
$\sfname{SquareSt.WeilConvergence}$.
\formallabel{GlobalTheory.WeilConv.twisted_weil_representation_convergence}
\formallabel{GlobalTheory.WeilConv.weilConvergence_of_siegelSet}
\end{theorem}

This is Weil's convergence criterion for theta integrals \cite{Weil1965} in the unitary case, with a polynomial bound in the
height. The hypothesis $F$ of the first formal name is not used. The proof rests on a bound for the theta kernel on Siegel
sets, which in turn rests on a bound at torus elements.

\begin{lemma}[Theta values at torus elements]\label{pf:sh-theta:torus}
In the situation of \cref{pf:sh-theta:weilconv}, let $(X_f,Y_f)$ be a Witt frame of $M$ of rank $r$, $c>0$, $K_0$ a compact
neighbourhood of $0$ in the symplectic space $(\eps M)\otimes X$, and $\mu$ a measure of covolume one with a fundamental domain for
its rational points, defining Sobolev seminorms $\mathrm{sob}_k$. There are $k,A\in\N$ and $C\ge0$ such that for every $n\ge1$, every
theta function $\Psi$ invariant under the Heisenberg translations by $nu$, $u\in K_0$ with zero archimedean part, every
$t\in\mathrm{cone}(c)$ and every torus element $a_t$,
\[
|(\omega(a_t,1)\Psi)(0)|\le C\,n^A\,\mathrm{sob}_k(\Psi)\prod_{i=1}^rt_i^{[F:\Q]\rank X}.
\]
\formallabel{GlobalTheory.T.thetaValue_siegelTorusElement_le_sobolevSeminorm}
\formallabel{GlobalTheory.SiegelTorusTheta.torus_bound}
\end{lemma}

\begin{proof}
Use the mixed model attached to the Witt decomposition of $\eps M$ defined by the frame. (1) There are a hermitian space $S_0$ (the
anisotropic kernel) and a hyperbolic decomposition $d$ of $\eps M$ such that every torus element $a_t$ is the Levi element $d(\alpha)$
of some $\alpha\in\GL_r(\A_E)$ with trivial finite part and archimedean components $\mathrm{diag}(t_1,\dots,t_r)$
(\lref{GlobalTheory.T.wittFrame_torus_eq_levi_element}). (2) For a theta function $f$ and a Levi element, with a suitable
covolume-one measure on $\mathrm{Mat}_{\rank X\times r}(\A_E)$,
\[
|(\omega(d(\alpha),1)f)(0)|\le|\alpha|^{\rank X}_{1/2}\sum_{\xi\in\mathrm{Mat}_{\rank X\times r}(E)}\bigl|\mathrm{mix}(f)(\xi\alpha)(0)\bigr|,
\]
where $\mathrm{mix}(f)$ is the mixed-model transform of $f$ attached to $d$ and $|\alpha|^{n}_{1/2}$ an explicit modulus factor
(\lref{GlobalTheory.T.leviThetaValue_le_det_mul_rational_sum}); and for $\alpha$ as in (1),
$|\alpha|^{\rank X}_{1/2}=\prod_it_i^{[F:\Q]\rank X}$ (\lref{GlobalTheory.T.halfPowerDetAbs_realDiagonal_torus}). (3) There
are $k,A,C$ such that for $\Psi,n$ as in the statement, $t\in\mathrm{cone}(c)$ and $\alpha$ as in (1), the sum in (2) is at most
$C\,n^A\,\mathrm{sob}_k(\Psi)$ (\lref{GlobalTheory.T.torus_mixedCoordinate_sum_le_sobolev_seminorm}). Combining (1)--(3) gives
the claim.
\end{proof}

\begin{lemma}[Theta kernels on Siegel sets]\label{pf:sh-theta:siegelbound}
In the situation of \cref{pf:sh-theta:weilconv}, let $(X_f,Y_f)$ be a Witt frame of $M$ of rank $r$, $\Omega,C_0$ compact,
$K$ a compact subgroup and $c>0$, with $a_t^{-1}pa_t\in C_0$ for $p\in\Omega$, $t\in\mathrm{cone}(c)$. Then there are $N$ and
$C$ such that for all $h$, $p\in\Omega$, $t\in\mathrm{cone}(c)$ and $k\in K$
\[
|\theta(pa_tk,h;\Phi)|\le C\,\shthetaht(h)^N\prod_{i=1}^rt_i^{[F:\Q]\rank X}.
\]
\formallabel{GlobalTheory.T.thetaKernel_siegelSet_le_torusPower_height}
\formallabel{GlobalTheory.SiegelThetaBound.theta_le_torus_height}
\end{lemma}

\begin{proof}
Put $C_1=C_0K$, a compact set. We use three results.
(a) The level of the translates grows polynomially: there are a compact neighbourhood $K_0$ of $0$ in the symplectic space and
$m\in\R$, $e\in\N$ such that for all $y\in C_1$ and $h$ there is an integer $1\le n\le m\,\shthetaht(h)^e$ for which
$\omega(y,h)\Phi$ is invariant under the Heisenberg translations by $nu$, $u\in K_0$ with zero archimedean part
(\lref{GlobalTheory.T.weilTheta_level_polynomial_in_height}).
(b) Fix a measure $\mu$ of covolume one with a fundamental domain for the rational points of the symplectic space
(\lref{GlobalTheory.Size.exists_isCovolOne}), and let $k,A,C_2$ be as in \cref{pf:sh-theta:torus} for $K_0$ and $\mu$.
(c) There are $N_2$ and $C_3\ge0$ with $\mathrm{sob}_k(\omega(y,h)\Phi)\le C_3\shthetaht(h)^{N_2}$ for $y\in C_1$
(\lref{GlobalTheory.T.sobolevSeminorm_weilTheta_polynomial_in_height}).

Let $h$, $p\in\Omega$, $t\in\mathrm{cone}(c)$ and $k\in K$, and put $y=a_t^{-1}pa_tk\in C_1$ and $\Psi=\omega(y,h)\Phi$, a theta
function. Since $(pa_tk,h)=(a_t,1)(y,h)$, $\theta(pa_tk,h;\Phi)=(\omega(a_t,1)\Psi)(0)$. With $n$ as in (a), (b) and (c) give
\[
|\theta(pa_tk,h;\Phi)|\le C_2(m\,\shthetaht(h)^e)^AC_3\shthetaht(h)^{N_2}\prod_it_i^{[F:\Q]\rank X},
\]
using $n\le|m|\shthetaht(h)^e$ and $\shthetaht(h)\ge1$. This is the claim with $N=eA+N_2$ and $C=C_2|m|^AC_3$.
\end{proof}

\begin{proof}[Proof of \cref{pf:sh-theta:weilconv}]
Let $\tau$ be a real place at which $M$ has signature $(\rank M-p,p)$. Choose a Witt frame $(X_f,Y_f)$ of $M$, of some rank
$r$ (Witt decomposition of hermitian spaces, \cite{Scharlau1985}; \lref{GlobalTheory.SiegelSet.exists_isWittFrame}); since an isotropic subspace has dimension at most the number of
negative eigenvalues at $\tau$, $r\le p$ (\lref{GlobalTheory.T.hermitianWittIndex_le_negative_signature}). Let
$K,\Omega,c,C_0$ and $\mathfrak S$ be as in $\sfname{SiegelSetUnitary}$ for this frame, $C_1$ the constant of its integral
bound for $\nu$, and $N,C$ as in \cref{pf:sh-theta:siegelbound}. Fix $h$. The function $g\mapsto\theta(g,h;\Phi)$ is
continuous (continuity of the Weil representation) and left invariant under $\Uni(M)(F)$ (automorphy of theta functions).
Since $\mathfrak S$ meets every $\Uni(M)(F)$-orbit, the integral of a non-negative invariant measurable function over the
fundamental domain is at most its integral over $\mathfrak S$
(\lref{GlobalTheory.T.lintegral_fundamentalDomain_le_of_meets_all_orbits}). Hence, with the bound of
\cref{pf:sh-theta:siegelbound} for $f=|\theta(\cdot,h;\Phi)|$ and $H(t)=\max(C,0)\shthetaht(h)^N\prod_it_i^{[F:\Q]\rank X}$,
\[
\int_{\shthetafd}|\theta(g,h;\Phi)|\,d\nu\le\int_{\mathfrak S}|\theta(g,h;\Phi)|\,d\nu
\le C_1\max(C,0)\,\shthetaht(h)^N\,I,
\]
where $I=\int_{\mathrm{cone}(c)}\prod_it_i^{[F:\Q]\rank X}\bigl(\delta(t)\prod_it_i\bigr)^{-1}dt$.
Since $\rank X<2(\rank M-p)\le2(\rank M-r)$, the integral $I$ is finite
(\lref{GlobalTheory.T.lintegral_siegelCone_inv_modulus_lt_top_weil}: on the cone the exponent of each $t_i$ in the
integrand is negative enough). This gives integrability and the bound with the constant $C_1\max(C,0)I$, independent of $h$.
\end{proof}

\subsubsection{Local matching for squares}

For the square statements, a \emph{triple} $(W,W_1,X)$ of hermitian spaces with unitary characters $(\alpha,\beta)$ satisfies
the rank, signature, infinity-type and ramification conditions of the formal definition (in particular
$\rank X=\rank W+\rank W_1$ is even), and \emph{real frames} $A,A_1,B$ of $W,W_1,X$ (signs and frames diagonalising the forms
at the real places) are fixed. For a sign $\eps$ and frames $A$ of $M$, $B$ of $X$, a \emph{finite vacuum function} $f$ is a
theta function which is a Fock vacuum at every real place (for the Cartan involutions of the frames), and for a polynomial
$P$ in the Fock variables $\mathrm{fock}(f,P)$ is the pure tensor of $f$ with the Fock vector of $P$; the span of these is the
\emph{algebraic Fock space} $\mathcal A_\eps(A,B)$. A \emph{matching datum} for $(\alpha,\beta)$ consists of Weil
representations $\omega_A$ of $(-W,X)$ with characters $(\beta,\beta/\alpha)$ and $\omega_B$ of $(W_1,X)$ with characters
$(\alpha,\alpha/\beta)$, an irreducible module $\Pi$ for $\Uni(X)(\A_{\mathrm f})$ and the Lie algebras of $\Uni(X)$ at the
real places (with actions $r_{\mathrm f}$, $r_{\mathrm l}$), and \emph{coinvariant maps} $p:\mathcal A_+(A,B)\to\Pi$,
$p':\mathcal A_-(A_1,B)\to\Pi$: surjective linear maps invariant under $\Uni(W)(\A_{\mathrm f})$, resp.\
$\Uni(W_1)(\A_{\mathrm f})$, killed by the derivatives along the Lie algebras of these groups, and equivariant for
$\Uni(X)(\A_{\mathrm f})$ and its Lie algebras. Thus $\Pi$ is a common irreducible quotient of the coinvariants of the two
algebraic Fock spaces.

\begin{definition}[Square local matching]\label{pf:sh-theta:sqlm}
$\sfname{SquareSt.LocalMatching}$ asserts: for every datum, standard $\psi$, archimedean data $\theta$, hermitian spaces
$W,W_1,X$ with real frames $A,A_1,B$ and functions $a,b$ from the real places to $\Q$, there are a vector space $\Pi_\R$ and
surjective linear maps $q$, $q'$ from the polynomial Fock spaces of $(W,X)$ and $(W_1,X)$ onto $\Pi_\R$ such that for all
$(\alpha,\beta)$ forming a triple with infinity exponents $a(\tau)$ and $b(\tau)$: matching data exist, and for every matching
datum there are a vector space $\Pi_{\mathrm f}$, linear maps $r,r'$ from the spaces of functions on the two symplectic
spaces to $\Pi_{\mathrm f}$, each surjective already on the finite vacuum functions, and an isomorphism
$e:\Pi_\R\otimes\Pi_{\mathrm f}\to\Pi$ with
\[
p(\mathrm{fock}(f,P))=e\bigl(q(P)\otimes r(f)\bigr),\qquad p'(\mathrm{fock}(f',P'))=e\bigl(q'(P')\otimes r'(f')\bigr)
\]
for all finite vacuum functions $f,f'$ and polynomials $P,P'$.
\formal{GlobalTheory.SquareSt.LocalMatching}\formal{GlobalTheory.SquareComm.IsMatch}\formal{GlobalTheory.SquareComm.IsTriple}
\end{definition}

The archimedean factor $(\Pi_\R,q,q')$ depends only on the spaces, frames and infinity exponents, not on the characters.
The proof is an algebraic decomposition of $\Pi$ into an archimedean and a finite part, which we prove first.

\begin{lemma}[Product form]\label{pf:sh-theta:productform}
Let $(\rho_i)_{i\in I}$ be endomorphisms of a complex vector space $\Pi_\R$ and $(\sigma_i)_{i\in I}$ endomorphisms of a complex
vector space $\Pi$. Assume that $\Pi_\R$ is irreducible (its only subspaces stable under all $\rho_i$ are $0$ and $\Pi_\R$) and
that every endomorphism commuting with all $\rho_i$ is a scalar. Call a linear map $\phi:\Pi_\R\to\Pi$ an \emph{intertwiner}
if $\phi\rho_i=\sigma_i\phi$ for all $i$. Let $q:Q\to\Pi_\R$ and $q':Q'\to\Pi_\R$ be surjective linear maps,
$\beta:V\to\Hom(Q,\Pi)$ and $\beta':V'\to\Hom(Q',\Pi)$ linear maps and $S\subseteq V$, $S'\subseteq V'$ subspaces such that
\begin{enumerate}[label=(\roman*)]
\item for every $f\in S$ there is an intertwiner $\phi$ with $\beta(f)=\phi\circ q$, and for every $f'\in S'$ an intertwiner
$\phi'$ with $\beta'(f')=\phi'\circ q'$;
\item the elements $\beta(f)(x)$, $f\in S$, $x\in Q$, span $\Pi$, and so do the elements $\beta'(f')(x')$.
\end{enumerate}
Then there are a vector space $\Pi_{\mathrm f}$, linear maps $r:V\to\Pi_{\mathrm f}$ and $r':V'\to\Pi_{\mathrm f}$ with
$r(S)=r'(S')=\Pi_{\mathrm f}$, and an isomorphism $e:\Pi_\R\otimes\Pi_{\mathrm f}\to\Pi$ such that
$\beta(f)(x)=e(q(x)\otimes r(f))$ for $f\in S$ and $\beta'(f')(x')=e(q'(x')\otimes r'(f'))$ for $f'\in S'$.
\formallabel{GlobalTheory.SquareLM.productForm}
\formallabel{GlobalTheory.ProductForm.eval_injective}
\formallabel{GlobalTheory.ProductForm.factorSpace_le}
\end{lemma}

\begin{proof}
\emph{Step 1: linearly independent intertwiners are independent over $\Pi_\R$.} Let $(\phi_k)_{k\in\kappa}$ be linearly
independent intertwiners and $(x_k)$ a finitely supported family in $\Pi_\R$ with $\sum_k\phi_k(x_k)=0$; we show $x=0$. Suppose
not, and choose such an $x\ne0$ with support $s$ of minimal size, and $k_0\in s$. Let $R$ be the space of families $d$ supported in
$s$ with $\sum_k\phi_k(d_k)=0$. It is stable under $d\mapsto(\rho_id_k)_k$, because the $\phi_k$ are intertwiners. The projection
$\pi:R\to\Pi_\R$, $d\mapsto d_{k_0}$, is injective: if $d_{k_0}=0$ then $d$ has smaller support, so $d=0$ by minimality. Its image
is stable under all $\rho_i$ and contains $x_{k_0}\ne0$, so it is $\Pi_\R$; thus $\pi$ is bijective. For $k\in s$ the map
$T_k=(d\mapsto d_k)\circ\pi^{-1}$ commutes with all $\rho_i$, hence $T_k=a_k\cdot1$, and $a_{k_0}=1$ since $T_{k_0}=1$. For every
$y\in\Pi_\R$ the family $\pi^{-1}(y)=(a_ky)_k$ lies in $R$, so $\bigl(\sum_ka_k\phi_k\bigr)(y)=0$. Hence $\sum_{k\in s}a_k\phi_k=0$
with $a_{k_0}=1$, contradicting linear independence.

\emph{Step 2: the evaluation map.} Let $F_S$ be the set of intertwiners $\phi$ with $\beta(f)=\phi\circ q$ for some $f\in S$; it is
a subspace ($\beta$ is linear), and similarly $F_{S'}$. For a subspace $F$ of intertwiners, the evaluation map
$\mathrm{ev}:\Pi_\R\otimes F\to\Pi$, $y\otimes\phi\mapsto\phi(y)$, is injective: writing an element as $\sum_ky_k\otimes b_k$ for a
basis $(b_k)$ of $F$, this is Step~1. If $F\supseteq F_S$, $\mathrm{ev}$ is surjective: $\beta(f)(x)=\phi(q(x))=\mathrm{ev}(q(x)\otimes\phi)$
for $f\in S$, and these elements span $\Pi$ by (ii).

\emph{Step 3: $F_S=F_{S'}$.} Let $\phi\in F_S$ and suppose $\phi\notin F_{S'}$. Then $\phi\ne0$; choose $y$ with $\phi(y)\ne0$. By Step 2
for $F_{S'}$ (surjectivity), $\phi(y)=\sum_kb_k(y_k)$ for a basis $(b_k)$ of $F_{S'}$ and some $y_k$. The family consisting of $\phi$
and the $b_k$ is linearly independent ($\phi\notin F_{S'}$), and $\phi(y)-\sum_kb_k(y_k)=0$; by Step 1, $y=0$, a contradiction. So
$F_S\subseteq F_{S'}$, and by symmetry $F_S=F_{S'}=:F$.

\emph{Step 4: the maps $r,r'$.} For $f\in S$ the intertwiner $\phi_f$ with $\beta(f)=\phi_f\circ q$ is unique because $q$ is
surjective; hence $f\mapsto\phi_f$ is a linear map $S\to F$, which we extend to a linear map $r:V\to F$. In the same way we obtain
$r':V'\to F$ with $\beta'(f')=r'(f')\circ q'$ for $f'\in S'$ (using $F_{S'}=F$). Put $\Pi_{\mathrm f}=F$ and $e=\mathrm{ev}$, an
isomorphism by Step 2. Then $\beta(f)(x)=r(f)(q(x))=e(q(x)\otimes r(f))$, and similarly for $\beta'$. Every $\phi\in F=F_S$ equals
$\phi_f=r(f)$ for some $f\in S$, and every $\phi\in F=F_{S'}$ equals $r'(f')$ for some $f'\in S'$.
\end{proof}

\begin{theorem}[Local matching for squares]\label{pf:sh-theta:sqlmthm}
Howe duality at the real place, $F$ and Theorem~K imply $\sfname{SquareSt.LocalMatching}$.
\formallabel{GlobalTheory.SquareLM.twisted_square_local_matching}
\end{theorem}

\begin{proof}
Fix the data $W,W_1,X,A,A_1,B,a,b$. Under the hypotheses, the \emph{real quotient} exists
(\lref{GlobalTheory.T.twisted_quotient_real}, about $3\,300$ Lean lines, from the Howe condition on the coinvariants of the
Fock spaces at the real places, a semisimplicity argument and the comparison of simple quotients): there are a vector space
$\Pi_\R$ of countable dimension with an action $\rho$ of the Lie algebras of $\Uni(X)$ at the real places under which it is
irreducible, and surjective linear maps $q,q'$ from the polynomial Fock spaces of $(W,X)$ and $(W_1,X)$ onto $\Pi_\R$, such that
for every triple $(\alpha,\beta)$ with infinity exponents $a,b$ and every matching datum, and every finite vacuum function $f$ of
the first kind (resp.\ $f'$ of the second kind), there is a linear map $\phi:\Pi_\R\to\Pi$ intertwining $\rho$ with $r_{\mathrm l}$
such that $p(\mathrm{fock}(f,P))=\phi(q(P))$ for all $P$ (resp.\ $p'(\mathrm{fock}(f',P'))=\phi(q'(P'))$). By Dixmier's lemma
(\lref{GlobalTheory.dixmier_commutant_scalar}: an endomorphism of a module of countable dimension over $\C$ which is
irreducible under a family of endomorphisms and commutes with them is a scalar; cf.\ \cite{Dixmier1974}) the commutant of $\rho$
consists of scalars. Matching data exist for every triple (\lref{GlobalTheory.T.twist_triple_matches}).

Given a matching datum, the finite vacuum functions form subspaces $S$, $S'$ and $\mathrm{fock}(f,P)$ is linear in $f$ on them
(the quasi-periodicity, level and Fock vacuum conditions are linear, and the pure tensor is linear in its finite component), so
$\beta(f)(P)=p(\mathrm{fock}(f,P))$ and $\beta'(f')(P')=p'(\mathrm{fock}(f',P'))$
are bilinear. Condition (i) of \cref{pf:sh-theta:productform} is the factorisation through the real quotient, and (ii) holds
because $p$ maps the algebraic Fock space, spanned by the $\mathrm{fock}(f,P)$, onto $\Pi$ (and similarly for $p'$).
\Cref{pf:sh-theta:productform}, with $I$ the set of pairs (real place, Lie algebra element) and $\sigma=r_{\mathrm l}$, gives
$\Pi_{\mathrm f}$, $r,r'$ on $S,S'$ and $e$; extending $r,r'$ linearly to all functions gives the statement.
\end{proof}

\begin{corollary}[Square statements]\label{pf:sh-theta:sqcor}
\begin{enumerate}[label=(\alph*)]
\item Howe duality at the real place, $F$ and Theorem~K imply $\sfname{SquareSt.LocalMatching}$.
\item $F$ and Theorem~K imply $\sfname{SquareSt.ThetaIntegralsAgree}$: for every triple, all Haar measures with fundamental
domains and every matching datum there is $c\ne0$ such that every $\Phi'$ in the algebraic Fock space of the second kind is
congruent, modulo the span of the $\omega_B(u,1)\Psi-\Psi$ ($u$ finite), to some $\Phi_0$ whose theta integral converges with
moderate growth and satisfies $\int\theta(g,\cdot;\Phi)\,dg=c\int\theta(g',\cdot;\Phi_0)\,dg'$ for every $\Phi$ of the first kind with
$p(\Phi)=p'(\Phi_0)$.
\item $F$ implies $\sfname{SumOKK}$ (\cref{def:pf:rowprops:five}).
\end{enumerate}
\formallabel{GlobalTheory.T.twisted_local_factors_match}
\formallabel{GlobalTheory.T.twisted_square_theta_integrals_eq}
\formallabel{GlobalTheory.T.sumOKK_of_theta_local_inputs}
\end{corollary}

\begin{proof}
(a) is \cref{pf:sh-theta:sqlmthm}. (b) is \lref{GlobalTheory.SquareAgree.twisted_theta_integrals_agree}, which states exactly
this. (c) is \lref{GlobalTheory.RowClaim.rowSum_of_hypotheses_alt}, which derives $\sfname{SumOKK}$ from $F$.
\end{proof}

\subsubsection{Cuspidality of the lifts along a grid}

Let $D$ be a datum, $R$ a realisation, $G$ a grid pre-datum over $R$ (\cref{def:pf:rowprops:grid}) and $\theta$ the archimedean
data. The nodes of $G$ carry discrete series models; for a node $N$ let $\mathrm{Cusp}(N)$ be its \emph{cuspidal multiplicity
space}: the vector-valued cusp forms $F$ which assign to each smooth vector $v$ of the model of $N$ an automorphic form $F(v)$ on
$\Uni(V_N)$. An arrow $a:N\to N'$ of the grid carries a Weil representation and, for every polynomial $P$, an intertwiner
$\mathrm{It}_a(P)$ from smooth vectors of the model of $N$ to those of $N'$. The arrow \emph{lifts cusp forms} for a finite vacuum
function $f$ if for every $F\in\mathrm{Cusp}(N)$ there is $F'\in\mathrm{Cusp}(N')$ such that for all $P$, all smooth vectors $v$ and all
$h$, the integral of $g\mapsto K_a(\mathrm{fock}(f,P);h,g)F(v)(g)$ over a fundamental domain converges and equals
$F'(\mathrm{It}_a(P)v)(h)$; for the upward arrows of the column $\mathrm c_i\to\mathrm c_{i+1}$ the kernel is the complex
conjugate of the kernel of the downward arrow.

\begin{definition}[Cuspidal lifts]\label{pf:sh-theta:liftscusp}
$\sfname{LiftsCuspidal}$ asserts: for every datum $D$ with $r>1800$ whose consecutive Harish-Chandra parameters differ by more
than $2N(D)^3$ at every real place, every realisation $R$ satisfying the row base conditions (level characters at the base type,
archimedean data, the measures of the models) and every grid pre-datum $G$ over $R$ satisfying its standing conditions, all
arrows of $G$ (horizontal $(i,j)\to(i+1,j)$, vertical $(i,j+1)\to(i,j)$, the two arrows between $\mathrm c_i$ and $\mathrm c_{i+1}$,
and $\mathrm t_j\to\mathrm t_{j+1}$, in range) lift cusp forms for every finite vacuum function.
\formal{GlobalTheory.IndefiniteRowSideFullHecke.CuspAlongSmoothLifts.LiftsCuspidal}
\formal{GlobalTheory.IndefiniteRowSideFullHecke.CuspAlongSmoothLifts.CuspAlong}
\end{definition}

\begin{theorem}[Theta lifts along a grid preserve cuspidality]\label{pf:sh-theta:cusp}
$F$ implies $\sfname{LiftsCuspidal}$.
\formallabel{GlobalTheory.PairLift.thetaPairLift_cuspidal_of_hypotheses}
\formallabel{GlobalTheory.T.twisted_thetaLift_cuspidal}
\formallabel{GlobalTheory.PairLift.liftsCuspidal_of_functoriality}
\formallabel{GlobalTheory.PairLift.liftsCuspidal}
\formallabel{GlobalTheory.PairLift.cuspAlong_of_pairLifts}
\end{theorem}

The mechanism is Rallis' tower principle \cite{Rallis1984}: the theta lift of a cusp form is cuspidal as soon as its lifts to all
lower members of the Witt tower vanish; here the vanishing of the lower lifts is first proved on Fock vectors, from the
archimedean inputs.

\begin{proof}
Only four components of $F$ are used: the discrete series cover, Harish-Chandra modules, rapid decay of cusp forms and
Casselman--Wallach globalisation. From the first two there is a family $X$ of real discrete series data (a ``chosen series'')
together with the archimedean Rallis properties of its theta correspondence
(\lref{GlobalTheory.T.exists_real_discreteSeries_data_theta_properties}), and from all four, for every datum $D$, the analytic
inputs of the pair lifts at $(D,X)$: the models are discrete series of the cover, vacuum characters, Lie and compact actions on
Fock spaces, density of Fock vectors, the rapid decay properties of the lifts, constant terms along towers and extensions of
$(\mathfrak g,K)$-modules
(\lref{GlobalTheory.T.thetaPairAnalyticHypotheses_of_discreteSeries}). With these inputs the following chain holds:
lower-tower lifts vanish on Fock vectors (\lref{GlobalTheory.T.lowerTowerLifts_vanish_on_fock}), hence on all theta
functions (\lref{GlobalTheory.T.lowerTowerLiftsVanish_of_vanish_on_fock}); hence theta lifts of cusp forms are cuspidal
(\lref{GlobalTheory.T.thetaLifts_cuspidal_of_lower_tower_lifts_vanish}); hence the \emph{pair lift statement} holds
(\lref{GlobalTheory.T.pairThetaLift_factors_through_intertwiner}): for spaces $V,W$ with discrete series models $M,M'$ whose
parameters are related by the theta rule of $X$ at all real places, a Weil representation with the corresponding infinity
exponents, a vacuum vector and the intertwiners $\mathrm{It}$ of the Fock vectors, the theta lift with any finite vacuum function
carries $\mathrm{Cusp}$ of $M'$ to $\mathrm{Cusp}$ of $M$ compatibly with $\mathrm{It}$.

It remains to apply the pair lift statement to the five kinds of arrows. For an arrow from $S$ to $T$ the parameters of the models
at the two ends, after passing to the contragredient where the arrow reverses the roles, are related by the theta rule of $X$ when
$D$ satisfies the size conditions (\lref{GlobalTheory.T.contragredientThetaParams_addBlockArrow}, \lref{GlobalTheory.T.contragredientThetaParams_removeBlockArrow}, \lref{GlobalTheory.T.contragredientThetaParams_columnArrows}, and \lref{GlobalTheory.T.contragredientThetaParams_lastRowArrows}).
The model of a negated space, or the complex conjugate model, is transported to a discrete series model with the same smooth
vectors (\lref{GlobalTheory.T.exists_negatedSpace_discreteSeriesModelTransport} and \lref{GlobalTheory.T.discrete_series_model_conjugate_transfer}),
and Fock vectors for frames with the same Cartan involutions correspond bijectively
(\lref{GlobalTheory.T.fockVector_frameChange_same_cartan}). For arrows lifting towards the source of the pair, this gives
the lift directly (\lref{GlobalTheory.PairLift.lifts_toSource}); for arrows towards the target, one passes to the conjugate
Weil datum of the exchanged pair and the adjoint Fock vectors (\lref{GlobalTheory.PairLift.lifts_toTarget}), and for the
upward column arrow the intertwiner is the adjoint of the real theta intertwiner of the downward arrow
(\lref{GlobalTheory.T.realThetaIntertwiner_adjoint_isRealThetaIntertwiner}). The infinity exponents of the characters
$\chi_{m-i}$, $\mu_t(j+1)$ and $\eta$ are those of the grid. This proves $\sfname{LiftsCuspidal}$; the second formal name of the
theorem restates it (\cref{pf:sh-theta:rem:wrappers}).
\end{proof}

% ---------------------------------------------------------------------------
\subsection{The Witt tower of a square}\label{sec:pf:sh-theta:tower}

This subsection proves the Witt-tower norm formula $\sfname{Tower}$ and the existence statement $\sfname{TriEx}$ for bottom
triples, which are items \ref{pf:fit:S:Tower} and \ref{pf:fit:S:TriEx} of the square statements of \cref{pf:fit:sqtop}
and are used in the proof of \cref{pf:rowprops:main}. The constant of a seesaw square of a grid compares two
theta integrals; it is computed by descending a Witt tower: hyperbolic planes are removed one at a time from the two spaces
of the square until totally definite spaces are reached, and each removal multiplies the constant by a Hecke $L$-value
and by an explicit arithmetic factor. The descent itself (\cref{pf:sh-theta:descent}) and the final computation are
proved here in full. The analytic and arithmetic inputs (the five statements of \cref{pf:sh-theta:five}) are reduced
to the seesaw identities, regularisation statements and local computations of the formal development, which are cited.

\subsubsection{Ranks, units and step multiples}

We use the grids of \cref{def:pf:rowprops:grid}, their Weil data (\cref{def:pf:rowprops:weil}), real matches and square
constants (\cref{def:pf:rowprops:square}) and the uniformity $\shthetaGBig$ of \cref{pf:fit:uniform}. Throughout, $D$ is a
global datum with ranks $k,m,t$ and $r=k+m-1$ (so $k,m$ are even, $t$ is odd, $m\ge2$, $m+3\le k$ and $3m<t$), $E$ is
its CM field and $\shthetaLL$ its coefficient field with the fixed embedding $\shthetaemb:\shthetaLL\to\C$; $T$ is the set
of types of the level characters, with base type $\mu_0$.

\begin{definition}[Ranks of a square]\label{pf:sh-theta:ranks}
Put $s=r+6$ and $r_{ij}=r+jt+is$ for $i,j\ge0$ (the rank of the node $(i,j)$ of a grid). For a square $(i,j)$,
$0\le i,j<m$, put
\[
p_j=\Bigl\lfloor\frac{m-j-1}2\Bigr\rfloor,\quad q_j=\Bigl\lfloor\frac{m-j}2\Bigr\rfloor,\quad
e_T=\frac{r_{i+1,j+1}-r_{ij}}2,\quad h=\frac{r_{i+1,j+1}+r_{ij}}2,
\]
$b=r_{i+1,j+1}-2p_j$, $b'=r_{ij}-2q_j$, $e_0=(b-b')/2$ and $h_0=(b+b')/2$.
\formal{GlobalTheory.ChiGrid.rankV}\formal{GlobalTheory.SqTop.eTop}\formal{GlobalTheory.FunObj.SqNum.halfRank}
\formal{GlobalTheory.FunObj.SqNum.botBig}\formal{GlobalTheory.FunObj.SqNum.botSmall}\formal{GlobalTheory.SqTop.e0}
\end{definition}

Since $r$, $s$ and $t$ are odd, $r_{i+1,j+1}-r_{ij}=t+s$ and $r_{i+1,j+1}+r_{ij}$ are even, so $e_T=(t+s)/2$ and
$h=r_{ij}+e_T$ are integers. As $m\le r\le r_{ij}$, one has $b+2p_j=r_{i+1,j+1}$, $b'+2q_j=r_{ij}$ and
$2h=r_{i+1,j+1}+r_{ij}$ (\lref{GlobalTheory.FunObj.SqNum.ranks_exact}). Moreover $p_j\le m<t\le e_T$ and
$e_T+q_j\le e_T+m\le e_T+r_{ij}=h$; we use these two inequalities below.

For a type $t$ let $\shthetaloc(t)$ be the finite set of places of $E$ attached to $t$ by the grid (the places over $S_0$
and the places of the twist $t$, \cref{def:pf:rowprops:square}), $\shthetatw(t)\subseteq\shthetaloc(t)$ the subset of
those not lying over $S_0$, and $Q(t)=\prod_{w\in\shthetatw(t)}\Nm w$. For $v\in S_0$ let $c(v)=1$ if $v=v_0$ and
$c(v)=2$ otherwise.

\begin{definition}[Away-units and step multiples]\label{pf:sh-theta:units}
Let $n\in\N$. An element $x\in\shthetaLL$ is an \emph{$n$-unit} if $x\,n^k=a$ for some $k\ge0$ and some
$a\in\OO_{\shthetaLL}$ dividing a power of $n$. For a type $t$ and $u,u_0\in\C$, $u$ is a \emph{step multiple} of $u_0$
(at $t$) if there are a $Q(t)$-unit $b$ and elements $g_v,g'_v\in\OO_{\shthetaLL}$ dividing $(\Nm v)^{c(v)}$ for
$v\in S_0$ such that
\[
u=\shthetaemb\Bigl(b\prod_{v\in S_0}\frac{g_v}{g'_v}\Bigr)\,u_0 .
\]
\formal{GlobalTheory.Square.IsAwayUnit}\formal{GlobalTheory.Square.StepRatio}\formal{GlobalTheory.SqTop.Qn}
\end{definition}

\begin{lemma}\label{pf:sh-theta:unitlemma}
Let $n\in\N$ and $t$ a type.
\begin{enumerate}[label=(\alph*)]
\item Products of $n$-units are $n$-units, and the inverse of a non-zero $n$-unit is an $n$-unit.
\item Every $u\in\C$ is a step multiple of itself. If $u$ is a step multiple of $u_0\ne0$, then $u/u_0$ is a step
multiple of $1$.
\item If $u_0$ is a step multiple of $u$, $u_0\ne0$, $C\in\C$ and $\zeta$ is a $Q(t)$-unit, then $C\,\shthetaemb(\zeta)\,u$
is a step multiple of $C\,u_0$.
\end{enumerate}
\formallabel{GlobalTheory.Square.IsAwayUnit.mul}
\formallabel{GlobalTheory.SqTower.isAwayUnit_inv}
\formallabel{GlobalTheory.SqTower.stepRatio_div}
\formallabel{GlobalTheory.SqStepBig.stepRatio_swap}
\end{lemma}

\begin{proof}
(a) If $xn^k=a$ and $x'n^{k'}=a'$ with $a\mid n^l$, $a'\mid n^{l'}$, then $xx'n^{k+k'}=aa'$ and $aa'\mid n^{l+l'}$. If
$x\ne0$, $xn^k=a$ and $n^l=ab$ with $b\in\OO_{\shthetaLL}$, then $x^{-1}n^l=n^kb$, and $n^kb$ divides $n^{k+l}$ because
$a\cdot n^kb=n^{k+l}$.
(b) Take $b=1$ and $g_v=g'_v=1$. If $u=\shthetaemb(b\prod g_v/g'_v)u_0$ with $u_0\ne0$, then
$u/u_0=\shthetaemb(b\prod g_v/g'_v)\cdot1$.
(c) Write $u_0=\shthetaemb(y)u$ with $y=b\prod_vg_v/g'_v$. As $u_0\ne0$, $y\ne0$, so $b\ne0$ and all $g_v\ne0$. Then
$C\shthetaemb(\zeta)u=\shthetaemb(\zeta b^{-1}\prod_vg'_v/g_v)\,Cu_0$, and $\zeta b^{-1}$ is a $Q(t)$-unit by (a);
the roles of $g_v$ and $g'_v$ are exchanged.
\end{proof}

\subsubsection{Theta triples, constants and factorisations}

Let $A$ be a hermitian space over $E$. A \emph{rigging} of $A$ consists of a lattice $\Lambda_A$, a frame
$\beta_A\in\GL_{\rank A}(\A_{E,\mathrm f})$, and at every real place $\tau$ a Cartan involution $s_{A,\tau}$ of the form
of $A$ with an orthonormal basis $B_{A,\tau}$; it is \emph{valid} if $\Lambda_A$ is self-dual with frame $\beta_A$ and
the Cartan data are compatible with the form, and \emph{standard} if the Gram matrix of $A$ is the identity,
$\Lambda_A=\OO_E^{\rank A}$ and $s_{A,\tau}=B_{A,\tau}=1$. For hermitian spaces $A,W$ let $\shthetaPol(W,A)$ be the
space of polynomials in the Fock variables of the pair $(A,W)$ (\cref{def:pf:rowprops:weil}).

\begin{definition}[Theta triples]\label{pf:sh-theta:triple}
Let $T$ be a set. A \emph{theta triple} $X$ over $T$ consists of hermitian spaces $A,A',W$ with riggings; for every
$t\in T$ representations $\omega_t$ of $\Uni(A)(\A)\times\Uni(W)(\A)$ and $\omega'_t$ of
$\Uni(A')(\A)\times\Uni(W)(\A)$ on spaces of functions on $\mathrm{Mat}_{\rank A\times\rank W}(\A_E)$, resp.\
$\mathrm{Mat}_{\rank A'\times\rank W}(\A_E)$; vacuum vectors $\mathrm{vac},\mathrm{vac}'$; an element $e_t$ of the
adeles, exponents $c_t(w)\in\N$ at the places $w$ of $E$ and finite Schwartz functions $f_t,f'_t$; and Haar measures
$m,m'$ with fundamental domains $\shthetafd,\shthetafd'$ for $\Uni(A)(F)$, $\Uni(A')(F)$. For $P\in\shthetaPol(W,A)$,
$P'\in\shthetaPol(W,A')$, $t\in T$ and $h\in\Uni(W)(\A)$ the \emph{theta kernels} are
\[
K_t(P;h)(g)=\theta\bigl(\omega_t(g,h)\,\Phi_t(P)\bigr),\qquad K'_t(P';h)(g')=\theta\bigl(\omega'_t(g',h)\,\Phi'_t(P')\bigr),
\]
where $\Phi_t(P)$ is the pure tensor of $f_t$ with the Fock vector of $P$ built from the archimedean riggings of $-A$ and
$W$, and $\Phi'_t(P')$ the analogous tensor for $(A',W)$; the \emph{theta integrals} are the functions
\[
I_t(P)(h)=\int_{\shthetafd}K_t(P;h)\,dm,\qquad I'_t(P')(h)=\int_{\shthetafd'}K'_t(P';h)\,dm'
\]
on $\Uni(W)(\A)$. Let $\alpha,\beta:T\to\{\text{unitary characters of }\A_E^\times\}$, $\nu:T\to\{\text{Hecke
characters of }E\}$ and $S:T\to\{\text{finite sets of places of }E\}$. The triple is \emph{of type
$(\alpha,\beta,\nu,S)$} if $\rank A+\rank A'=\rank W$; the three riggings are valid; for every $t$, $\omega_t$ is a Weil
representation of the pair $(-A,W)$ with splitting characters $(\beta_t,\beta_t/\alpha_t)$ and $\omega'_t$ one of
$(A',W)$ with splitting characters $(\alpha_t,\alpha_t/\beta_t)$; the vacuum vectors are vacuum vectors of the
riggings; $e_t$ is an idempotent sum for the places of $S(t)$; $c_t(w)$ is the conductor exponent of $\nu(t)$ at $w\in
S(t)$; $f_t$ and $f'_t$ are the theta functions at $e_t$ of the vacuum vectors and of the two explicit unit-character
sections attached to $\alpha_t/\beta_t$ and the ideal $\prod_{w\in S(t)}w^{c_t(w)}$; and $m,m'$ are Haar measures with
fundamental domains.
\formal{GlobalTheory.IndefiniteRowSideFullHecke.UnitCharTriangle.Tri}
\formal{GlobalTheory.IndefiniteRowSideFullHecke.UnitCharTriangle.Tri.Is}
\formal{GlobalTheory.IndefiniteRowSideFullHecke.UnitCharTriangle.Rig.Is}
\end{definition}

\begin{definition}[Constants and factorisations]\label{pf:sh-theta:const}
Let $X$ be a theta triple over $T$ (with fixed $\psi,\theta$).
\begin{enumerate}[label=(\alph*)]
\item A function $c:T\to\C$ is a \emph{constant} of $X$ if $c$ is not identically zero and there are
$P\in\shthetaPol(W,A)$, $P'\in\shthetaPol(W,A')$ such that for every $t$: the kernels $K_t(P;h)$ and $K'_t(P';h)$ are
integrable on the fundamental domains for every $h$, $I'_t(P')\ne0$, and $I_t(P)=c(t)\,I'_t(P')$.
\item A \emph{factorisation} of $X$ with multiplier $\lambda:T\to\C$ consists of a complex vector space $\Pi$, linear maps
$p:\shthetaPol(W,A)\to\Pi$ (surjective) and $p':\shthetaPol(W,A')\to\Pi$, and injective linear maps $J_t$ from $\Pi$ to
the functions on $\Uni(W)(\A)$, such that $\lambda(t)\ne0$ for all $t$; all kernels $K_t(P;h)$ are integrable;
$I_t(P)=\lambda(t)J_t(p(P))$ for all $t,P$; $I'_t(P')=J_t(p'(P'))$ whenever the kernels $K'_t(P';h)$ are integrable for all
$h$; and there is $P'$ with $p'(P')\ne0$ whose kernels $K'_t(P';h)$ are integrable for all $t,h$.
\end{enumerate}
\formal{GlobalTheory.IndefiniteRowSideFullHecke.UnitCharTriangle.Tri.IsConst}\formal{GlobalTheory.ConstRatio.Factor}
\end{definition}

So a factorisation says that the theta integrals of both members of the triple factor through one common quotient
$\Pi$ of the space of polynomials, on which the integrals of the second member are realised by the same injective maps
$J_t$ as those of the first member, up to the scalars $\lambda(t)$.

\begin{lemma}[Constants of a factorised triple]\label{pf:sh-theta:factor}
Let $X$ be a theta triple with a factorisation with multiplier $\lambda$.
\begin{enumerate}[label=(\alph*)]
\item If $T$ is non-empty, $\lambda$ is a constant of $X$.
\item If $c$ is a constant of $X$, there is $\kappa\ne0$ with $c(t)=\kappa\lambda(t)$ for all $t$.
\end{enumerate}
\formallabel{GlobalTheory.ConstRatio.Factor.isConst}
\formallabel{GlobalTheory.ConstRatio.Factor.eq_mul}
\end{lemma}

\begin{proof}
(a) Let $P'$ be as in the last condition of the factorisation and choose $P$ with $p(P)=p'(P')$ ($p$ is surjective). For
every $t$ the kernels of $P$ and $P'$ are integrable, and
$I_t(P)=\lambda(t)J_t(p(P))=\lambda(t)J_t(p'(P'))=\lambda(t)I'_t(P')$. If $I'_t(P')=0$, then $J_t(p'(P'))=0$, so $p'(P')=0$
by injectivity of $J_t$, a contradiction; hence $I'_t(P')\ne0$. Finally $\lambda$ does not vanish, and $T\ne\emptyset$.

(b) Let $P,P'$ be the polynomials of the constant $c$ and $t_1$ a type with $c(t_1)\ne0$. For every $t$,
$J_t(\lambda(t)p(P))=I_t(P)=c(t)I'_t(P')=J_t(c(t)p'(P'))$, the kernels of $P'$ being integrable; by injectivity
$\lambda(t)p(P)=c(t)p'(P')$. Moreover $p'(P')\ne0$, since otherwise $I'_{t_1}(P')=J_{t_1}(0)=0$. Hence
$p(P)=\lambda(t)^{-1}c(t)\,p'(P')$ for every $t$, and comparing with $t=t_1$ gives
$\lambda(t)^{-1}c(t)=\lambda(t_1)^{-1}c(t_1)=:\kappa$, which is non-zero.
\end{proof}

\begin{definition}[Stage triples and bottom triples]\label{pf:sh-theta:stage}
Let $G$ be a grid datum with Weil datum $W$ (with character $\psi$ and archimedean data $\theta$), $(i,j)$ a square
and $p,q\in\N$. Put $\eps=r_{ij}\bmod 2$,
\[
\alpha_t=\mu_t(j+1)\,\eta^{\eps},\qquad\beta=\chi_{m-i}\,\eta^{\eps},\qquad
\nu_H(t,i,j)=\mu_{0,j}\,\xi_t\,\chi_{m-i}^{-1},
\]
where $\mu_t(j+1)$, $\eta$, $\xi_t$ and $\mu_{0,j}$ are the characters of the level characters and $\chi_{m-i}$ that of
the grid. A theta triple $X$ with spaces $(A,A',W')$ is a \emph{stage-$(p,q)$ triple} of the square if $p\le p_j$,
$q\le q_j$,
\[
\rank A+2p_j=r_{i+1,j+1}+2p,\qquad\rank A'+2q_j=r_{ij}+2q,
\]
the forms of $A$, $A'$ and $W'$ have exactly $p$, $q$ and $p+q$ negative eigenvalues at every real place, and $X$ is of
type $(\alpha,\beta,\nu_H(\cdot,i,j),\shthetaloc)$. It is a \emph{bottom triple} if $\rank A=b$, $\rank A'=b'$, its three
riggings are standard, and it is of the same type.
\formal{GlobalTheory.SqTower.IsStageC}\formal{GlobalTheory.ConstRatio.NegIdx}\formal{GlobalTheory.RowTest.IsBottomTri}
\formal{GlobalTheory.FunObj.SqNum.nuH}
\end{definition}

Stage $(p_j,q_j)$ has the ranks $r_{i+1,j+1}$ and $r_{ij}$ of the nodes of the square, and passing from stage
$(p,q+1)$ to $(p,q)$, or from $(p+1,0)$ to $(p,0)$, lowers the rank of $A'$, resp.\ of $A$, and that of $W'$ by $2$: one
hyperbolic plane is removed.

\begin{lemma}\label{pf:sh-theta:botstage}
Every bottom triple of a square is a stage-$(0,0)$ triple of it.
\formallabel{GlobalTheory.SqTower.IsBottomTri.isStageC}
\end{lemma}

\begin{proof}
The Gram matrices of its spaces are the identity, so the forms are positive definite at every real place: $0$, $0$ and $0$
negative eigenvalues. Its ranks are $b$ and $b'$, and $b+2p_j=r_{i+1,j+1}$, $b'+2q_j=r_{ij}$; and $0\le p_j$, $0\le q_j$. The type
is the same.
\end{proof}

\begin{definition}[Regularised constants; tied real matches]\label{pf:sh-theta:reg}
In the notation of \cref{def:pf:rowprops:square}, a number $c_0$ is a \emph{regularised constant} of the square $(i,j)$
at the type $t$ for a real match $\varrho$ if for every volume datum $M$ of level $t$, every matched pair
$(\Theta,\Theta')$ and all $g,g_1$ there is a finite combination $z$ with $\shthetadeg z\ne0$ such that the two route
kernels are integrable and
$\shthetadeg z\int K^{\mathrm{big}}_\Theta(g,g_1)\,d\nu=c_0\int z\cdot K^{\mathrm{sm}}_{\Theta'}(g,g_1)\,d\nu$. The real
match is \emph{tied} if for all squares, types and grid frames, and for every local matching datum of the square
(\cref{def:families:K}), the images in the common quotient of the finite theta data of matched pairs of route tensors
are proportional, with a factor independent of the pair.
\formal{GlobalTheory.SqTower.RegConst}\formal{GlobalTheory.FunObj.RealMatch.Tied}
\end{definition}

\subsubsection{The five tower statements}

In the following statements ``at large scale'' means: the statement $\shthetaGBig(Q)$ of \cref{pf:fit:uniform} holds for
the property $Q(G,\mu_0,C,x,b)$ which asserts the displayed conclusion for every Weil datum $W$ of $G$ with standard
character $\psi$ and every grid choice $S$ satisfying the field condition $\sfname{NumbersIn}$ at $\mu_0$ (the values of
the characters of the grid, the eighth roots of unity and the roots of unity of order $\Nm w^{\mathrm{depth}(w)}$,
$w\in\shthetaloc(t)$, lie in $\shthetaemb(\shthetaLL)$), and every square $(i,j)$. A finite conjunction of statements at
large scale holds at large scale (take the maxima of the thresholds), and a statement implied by one at large scale holds
at large scale. We write $F$ for
$\sfname{ThetaLocalInputs}$ (\cref{sec:families:theta}) and ``Theorem~K'' for the conjunction of
$\sfname{ThetaIntegralsAgree}$ and $\sfname{LocalMatching}$ (\cref{def:families:K}). For a Hecke character $\nu$, an
\emph{$L$-function of $\nu$} is an entire function which equals the Euler product of the Hecke $L$-function of $\nu$ on
$\Re s>1$; it is unique by the identity theorem. $L^{S_0}(\nu,s)$ denotes the partial $L$-function without the Euler
factors at $S_0$, and $L_t$ an $L$-function of $\nu_H(t,i,j)$.

\begin{proposition}[The five tower statements]\label{pf:sh-theta:five}
Assume $F$ and Theorem~K. Then the following hold at large scale.
\begin{enumerate}[label=\textup{(T\arabic*)},leftmargin=*]
\item\label{pf:sh-theta:T1} \emph{(Top stage.)} For every real match $\varrho$ there are a stage-$(p_j,q_j)$ triple $X$, a
constant $c$ of $X$ and $\kappa:T\to\C$ such that for every $t$, $\kappa(t)\ne0$ and $\kappa(t)=\shthetaemb(\zeta_t)\kappa(\mu_0)$
for a $Q(t)$-unit $\zeta_t$, and every square constant $c'$ of a non-zero square at $(i,j,t)$ equals $\kappa(t)c(t)$.
\item\label{pf:sh-theta:T2} \emph{(Factorisation.)} Every stage-$(p,q)$ triple has a factorisation.
\item\label{pf:sh-theta:T3} \emph{(Small-rank step.)} Let $X$ be a stage-$(p,q+1)$ triple with a factorisation with multiplier
$\lambda$, with $\rank A=\rank A'+2e$, and let $L_t$ be $L$-functions of $\nu_H(t,i,j)$. Then there are a stage-$(p,q)$
triple $Y$ with a constant $c$ and $y:T\to\C$ such that for all $t$: $y(t)\ne0$,
$\lambda(t)L_t(e+1)=y(t)c(t)$, and $y(t)$ is a step multiple of $y(\mu_0)$.
\item\label{pf:sh-theta:T4} \emph{(Large-rank step.)} Let $X$ be a stage-$(p+1,0)$ triple with a factorisation with multiplier
$\lambda$, with $\rank A=\rank A'+2e$, and let $L_t$ be as in (T3). Then there are a stage-$(p,0)$ triple $Y$ with a
constant $c$ and $y:T\to\C$ such that for all $t$: $y(t)\ne0$, $\lambda(t)=y(t)L_t(e)c(t)$, and $y(t)$ is a step multiple
of $y(\mu_0)$.
\item\label{pf:sh-theta:T5} \emph{(Bottom.)} If $X$ is a bottom triple and $Y$ a stage-$(0,0)$ triple with a constant $c$,
there is $v:T\to\C$ such that $t\mapsto v(t)c(t)$ is a constant of $X$, and for every $t$, $v(t)\ne0$ and
$v(t)=\shthetaemb(\zeta_t)v(\mu_0)$ for a $Q(t)$-unit $\zeta_t$.
\end{enumerate}
\formallabel{GlobalTheory.T.witt_tower_top_stage_theta_norm_formula}
\formallabel{GlobalTheory.SqTower.theta_tower_top_stage_norm_eq}
\formallabel{GlobalTheory.T.exists_factor_of_isStageC}
\formallabel{GlobalTheory.StageFactor.theta_tower_stage_norm_factorization}
\formallabel{GlobalTheory.T.witt_tower_norm_step_of_small_rank}
\formallabel{GlobalTheory.SqStep.theta_tower_norm_step_small_rank}
\formallabel{GlobalTheory.T.witt_tower_norm_step_of_large_rank}
\formallabel{GlobalTheory.SqStepBig.theta_tower_norm_step_large_rank}
\formallabel{GlobalTheory.T.witt_tower_constant_ratio_eq_bottom}
\formallabel{GlobalTheory.SqTower.theta_tower_bottom_constant_ratio}
\end{proposition}

The five statements rest on the following results of the formal development, which are cited (all of them hold at large
scale under $F$ and Theorem~K, except where stated otherwise).
\begin{enumerate}[label=(\roman*),leftmargin=*]
\item\label{pf:sh-theta:c:grid} For every Weil datum there is a stage-$(p_j,q_j)$ triple on the spaces
$(V_{i+1,j+1},V_{ij},V_{i,j+1}\oplus V_{i+1,j})$ of the square whose riggings are the lattices and Cartan involutions of
the grid (\lref{GlobalTheory.T.exists_grid_of_square_tower}; this needs only $F$).
\item\label{pf:sh-theta:c:tied} For every real match $\varrho$ there is a tied real match $\varrho_0$ with the same vacuum
vectors, and every square has a regularised constant for $\varrho_0$ at every type
(\lref{GlobalTheory.T.exists_tied_theta_pair}).
\item\label{pf:sh-theta:c:real} If $\varrho,\varrho_0$ have the same vacuum vectors, there is $l\ne0$ such that for every type
$t$, every square constant $c'$ for $\varrho$ and every regularised constant $c_0$ for $\varrho_0$ at $(i,j,t)$ one has
$c_0\ne0$ and $c'=l\,c_0$ (\lref{GlobalTheory.T.theta_real_places_separation}; separation at the real places).
\item\label{pf:sh-theta:c:reg} If $\varrho$ is tied and $X$ is the triple of \ref{pf:sh-theta:c:grid} with a factorisation with
multiplier $\lambda$, there is $\kappa$ with $\kappa(t)\ne0$, $\kappa(t)=\shthetaemb(\zeta_t)\kappa(\mu_0)$ for $Q(t)$-units
$\zeta_t$, such that every regularised constant at $(i,j,t)$ for $\varrho$ equals $\kappa(t)\lambda(t)$
(\lref{GlobalTheory.T.regConst_eq_away_unit_mul_factor}; the comparison of regularised constants, about $6\,000$ Lean
lines).
\item\label{pf:sh-theta:c:stage} \emph{(Stage dictionary.)} For every stage triple $X$ there are real frames of $A,A',W'$ and
functions $a,b$ on the real places such that for every $t$: $(A,A',W')$ with the characters $(\alpha_t,\beta)$ is a triple in the
sense of \cref{sec:pf:sh-theta:squares} with infinity exponents $a,b$; $f_t$ and $f'_t$ are finite vacuum functions; and
$K_t(P;h)(g)=\theta(g,h;\mathrm{fock}(f_t,P))$, $K'_t(P';h)(g')=\theta(g',h;\mathrm{fock}(f'_t,P'))$
(\lref{GlobalTheory.T.towerStage_thetaDictionary_exists} and \lref{GlobalTheory.T.towerStage_seesawDictionary_tripleConditions}).
Moreover there is $P'$ whose kernels $K'_t(P';h)$ are integrable for all $t,h$ with $I'_t(P')\ne0$ for all $t$
(\lref{GlobalTheory.T.theta_small_rank_stage_holds}; non-vanishing of the small-rank integrals, about $4\,700$ lines).
\item\label{pf:sh-theta:c:ktype} \emph{(Factorisation through the archimedean quotient.)} For a triple with real frames and infinity
exponents $a,b$ there are a vector space $\Pi_\R$ and a linear map $q$ from the polynomials to $\Pi_\R$ such that for every Weil
representation of the first pair with the characters of the triple and every finite vacuum function $f$ there are a vector space
$\Pi_{\mathrm f}$, $y\in\Pi_{\mathrm f}$ and a linear map $J$ from $\Pi_\R\otimes\Pi_{\mathrm f}$ to the functions on $\Uni(W')(\A)$,
zero or injective, with $\int\theta(g,\cdot;\mathrm{fock}(f,P))\,dg=J(q(P)\otimes y)$ for all $P$
(\lref{GlobalTheory.T.thetaIntegral_factors_through_kType}).
\item\label{pf:sh-theta:c:match} \emph{(Matching data.)} For a triple and Weil representations of its two pairs with the
characters of the triple, matching data with these Weil representations exist
(\lref{GlobalTheory.T.triple_weil_transfer_matching}). For a stage triple, a matching datum at $t$ whose coinvariant maps are
product maps as in \cref{pf:sh-theta:sqlm} satisfies $r(f_t)=\kappa\,r'(f'_t)$ for some $\kappa\ne0$
(\lref{GlobalTheory.T.theta_finiteData_match_each_stage}). For every matching datum there is $c\ne0$ such that
$\int\theta(g,\cdot;\Phi)\,dg=c\int\theta(g',\cdot;\Phi')\,dg'$ whenever $\Phi$, $\Phi'$ lie in the algebraic Fock spaces, the kernels
of $\Phi'$ are integrable and $p(\Phi)=p'(\Phi')$ (\lref{GlobalTheory.T.thetaIntegrals_agree_of_convergence}).
\item\label{pf:sh-theta:c:seesaw} \emph{(Seesaw reduction.)} In the situation of (T3) there are a stage-$(p,q)$ triple $Y$ with a
constant $c$ which vanishes nowhere and, for every $t$, a base datum fitting $(G,W,i,j,t)$, units $e_1,e_2,s\in\Z^\times$ and
fractional ideals $\mathfrak b,\mathfrak c$ of $E$ such that the step factors $\phi_t=\mathrm{stepFac}(t;-e)$ and
$\phi_{\mu_0}=\mathrm{stepFac}(\mu_0;-e)$ are non-zero, $L^{S_0}(\nu_t,e+1)\ne0\ne L^{S_0}(\nu_{\mu_0},e+1)$, and
\[
\lambda(t)L^{S_0}(\nu_t,e+1)\,\phi_t\,c(\mu_0)=\lambda(\mu_0)L^{S_0}(\nu_{\mu_0},e+1)\,\phi_{\mu_0}\,c(t)
\]
(\lref{GlobalTheory.T.seesaw_first_reduction_level}). Here $\nu_t=\alpha_t/\beta$, and the \emph{step factor} of the
step relations of Theorem~K (\cref{def:families:K}) is
$\mathrm{stepFac}(t;z)=\nu_t(\mathfrak b)^{e_1}\nu_t(\mathfrak c)^{e_2}\prod_{v}\nu_{t,v}(s)\,\epsilon(\psi_v,\nu_{t,v},z)$,
the product over the places $v$ of $F$ in $S_0$ or below a twist place of $t$, $\nu_{t,v}$ being the local component and
$\epsilon$ the local epsilon factor.
\item\label{pf:sh-theta:c:big} \emph{(Step identity.)} In the situation of (T4) there are a stage-$(p,0)$ triple $Y$ with a
constant $c$ and, for every $t$, base data, units, ideals, $C\in\C$ and a $Q(t)$-unit $\zeta$ with
$\lambda(t)=C\,\shthetaemb(\zeta)\,\phi_t\,L^{S_0}(\nu_t,e)\,c(t)$ and
$\lambda(\mu_0)=C\,\phi_{\mu_0}\,L^{S_0}(\nu_{\mu_0},e)\,c(\mu_0)$, where now $\phi_t=\mathrm{stepFac}(t;1-e)$
(\lref{GlobalTheory.SqStepBig.theta_tower_step_norm_identity}).
\item\label{pf:sh-theta:c:gauss} \emph{(Arithmetic of the step factors.)} If a base datum fits $(G,W,i,j,t)$ and the step
factor $\phi_t$ (for any shift) is non-zero, then $\phi_{\mu_0}$ is a step multiple of $\phi_t$
(\lref{GlobalTheory.T.thetaStepFactor_ratio_unit_gaussSumQuotient}; a quotient of Gauss sums is a unit away from
$Q(t)$; this holds at large scale without hypotheses).
\item\label{pf:sh-theta:c:bottom} \emph{(Bottom transport.)} For a bottom triple $X$ with spaces $(A,A',W')$ and a stage-$(0,0)$
triple $Y$: there is a stage-$(0,0)$ triple $Y'$ on the spaces of $X$ having every constant of $Y$
(\lref{GlobalTheory.T.definiteBottom_stage_transport_same_constants}); if $Y$ is on the spaces of $X$, there are a
stage-$(0,0)$ triple $Z$ with the lattices of $X$ and $v_1$ with $v_1(t)\ne0$, $v_1(t)=\shthetaemb(\zeta_t)v_1(\mu_0)$ for
$Q(t)$-units $\zeta_t$, such that $v_1c$ is a constant of $Z$ for every constant $c$ of $Y$
(\lref{GlobalTheory.T.definiteBottom_lattice_change_constants_awayUnit}); if moreover $Y$ has the lattices of $X$, the
same holds with $Z$ having the lattices and the frames of $X$
(\lref{GlobalTheory.T.definiteBottom_framesChange_stageZero}); and if $Y$ has the lattices and frames of $X$, then for
every constant $c$ of $Y$ there is $\kappa\ne0$ such that $\kappa c$ is a constant of $X$
(\lref{GlobalTheory.T.definiteBottom_tower_constants_differ_by_factor}). These four statements hold for all data, without
hypotheses.
\end{enumerate}

The partial $L$-values occurring in \ref{pf:sh-theta:c:seesaw} and \ref{pf:sh-theta:c:big} are values of $L$-functions of
$\nu_H(t,i,j)$:

\begin{lemma}[Partial $L$-values of a square]\label{pf:sh-theta:partialL}
The following holds at large scale, with all thresholds equal to $0$ and without hypotheses: for every type $t$, every integer
$s\ge2$ and every $L$-function $\Lambda$ of $\nu_H(t,i,j)$, one has $L^{S_0}(\nu_t,s)=\Lambda(s)$, where $\nu_t=\alpha_t/\beta$.
\formallabel{GlobalTheory.T.partialHeckeLFunction_square_eq_primitive}
\formallabel{GlobalTheory.SqStep.charL_chs_eq}
\end{lemma}

\begin{proof}
As $\alpha_t$ and $\beta$ contain the same power of $\eta$, $\nu_t=\mu_t(j+1)\,\chi_{m-i}^{-1}$. By the defining properties of the
level characters at the base type and of the characters of the grid, this idele class character and the Hecke character
$\nu_H(t,i,j)$ have the same local components (\lref{GlobalTheory.FunObj.SqNum.chrEq_nu}). Consequently $L^{S_0}(\nu_t,\cdot)$, the product of the local factors
$(1-\nu_t(\varpi_w)\Nm w^{-s})^{-1}$ at the places $w$ of $E$ not above $S_0$ at which $\nu_t$ is unramified, is the Hecke
$L$-function of $\nu_H(t,i,j)$ with the Euler factors at the places above $S_0$ removed
(\lref{GlobalTheory.LFun.charL_eq_of_chrEq}). The Hecke character $\nu_H(t,i,j)$ is ramified at every place of $E$ above a
place of $S_0$ (\lref{GlobalTheory.SqStep.nuH_not_unrAt_above}), so its Euler factors at these places are equal to $1$, and
removing them does not change the product. Since $\Re s=s>1$, the full Euler product converges unconditionally to $\Lambda(s)$,
by the definition of an $L$-function of $\nu_H(t,i,j)$; hence so does the partial one
(\lref{GlobalTheory.SqStep.heckeL_eq_of_ram}: if a Hecke character is ramified at every place of a set $S$, its
$L$-function without the factors at $S$ equals its $L$-function on $\Re s>1$).
\end{proof}

\begin{lemma}[Theta integrals of stage triples]\label{pf:sh-theta:stageint}
Assume $F$ and Theorem~K. At large scale, every stage-$(p,q)$ triple $X$ satisfies:
\begin{enumerate}[label=(\alph*)]
\item all kernels $K_t(P;h)$ are integrable on the fundamental domain;
\item there are a linear map $\sigma$ from the polynomials of $(A',W')$ to those of $(A,W')$ and $\lambda:T\to\C^\times$ with
$I_t(\sigma P')=\lambda(t)I'_t(P')$ whenever the kernels $K'_t(P';h)$ are integrable for all $h$;
\item if $I_t(P)=0$ for some $t$, then $I_{t'}(P)=0$ for all $t'$.
\end{enumerate}
\formallabel{GlobalTheory.T.theta_integral_stage_convergent}
\formallabel{GlobalTheory.StageConv.theta_stage_converges}
\formallabel{GlobalTheory.T.theta_stage_liftable}
\formallabel{GlobalTheory.StageLift.theta_stage_lift_exists}
\formallabel{GlobalTheory.T.theta_stage_kernel_characterization}
\formallabel{GlobalTheory.StageKer.theta_stage_kernel_eq}
\end{lemma}

\begin{proof}
(a) The bundle $F$ contains $\sfname{SiegelSetUnitary}$, so \cref{pf:sh-theta:weilconv} applies to the Weil representation
$\omega_t$ of $(-A,W')$, to the theta function $\Phi_t(P)$ (a pure tensor of the theta function $f_t$ with a Fock vector, hence a
theta function), and to the number $p$ of negative eigenvalues of $A$, provided $\rank W'<2(\rank A-p)$. By the stage conditions
and $\rank W'=\rank A+\rank A'$,
\[
2(\rank A-p)-\rank W'=\rank A-2p-\rank A'=t+s-2p_j+2(q_j-q)\ge t+s-(m-1)>0,
\]
since $2p_j\le m-1$, $q\le q_j$ and $s=r+6>m$.

(b) Let the frames and exponents of the dictionary \ref{pf:sh-theta:c:stage} be fixed, and let $(\Pi_\R,q,q')$ be the archimedean
factor of the square local matching (\cref{pf:sh-theta:sqcor}(a)) for these data; $q$ is surjective, so it has a linear right
inverse $g$. Put $\sigma=g\circ q'$. Fix $t$. By \ref{pf:sh-theta:c:match} there is a matching datum with the Weil representations
$\omega_t,\omega'_t$ of $X$, and by the square local matching its coinvariant maps are product maps $p(\mathrm{fock}(f,P))=
e(q(P)\otimes r(f))$, $p'(\mathrm{fock}(f',P'))=e(q'(P')\otimes r'(f'))$; by \ref{pf:sh-theta:c:match} again $r(f_t)=\kappa\,r'(f'_t)$ with
$\kappa\ne0$, and there is $c\ne0$ with the agreement property. Then, as $q(g(q'(P')))=q'(P')$,
\[
p\bigl(\mathrm{fock}(f_t,\sigma P')\bigr)=e\bigl(q'(P')\otimes\kappa\,r'(f'_t)\bigr)=p'\bigl(\kappa\,\mathrm{fock}(f'_t,P')\bigr).
\]
If the kernels of $P'$ are integrable, so are those of $\kappa\,\mathrm{fock}(f'_t,P')$, and the agreement property gives
$I_t(\sigma P')=c\kappa\,I'_t(P')$, using the dictionary to identify the kernels. Put $\lambda(t)=c\kappa\ne0$.

(c) Let $(\Pi_\R,q)$ be as in \ref{pf:sh-theta:c:ktype} for the triple at the frames of the dictionary; for each $t$ we get
$\Pi_{\mathrm f}$, $y$ and $J$ with $I_t(P)=J(q(P)\otimes y)$ for all $P$. Let $P'_0$ be the polynomial of \ref{pf:sh-theta:c:stage}
and $\sigma,\lambda$ as in (b); then $I_t(\sigma P'_0)=\lambda(t)I'_t(P'_0)\ne0$. Hence $J\ne0$, so $J$ is injective, and $y\ne0$.
Therefore $I_t(P)=0$ if and only if $q(P)\otimes y=0$, if and only if $q(P)=0$, a condition independent of $t$.
\end{proof}

\begin{lemma}[A factorisation from its parts]\label{pf:sh-theta:parts}
Let $X$ be a theta triple over $T$ and assume: all kernels $K_t(P;h)$ are integrable; if $I_t(P)=0$ for some $t$ then
$I_{t'}(P)=0$ for all $t'$; there are a linear map $\sigma:\shthetaPol(W,A')\to\shthetaPol(W,A)$ and $\lambda:T\to\C^\times$ with
$I_t(\sigma P')=\lambda(t)I'_t(P')$ whenever the kernels $K'_t(P';h)$ are integrable for all $h$; and there is $P'_0$ whose kernels
are integrable for all $t,h$ with $I'_t(P'_0)\ne0$ for all $t$. Then $X$ has a factorisation with multiplier $\lambda$.
\formallabel{GlobalTheory.StageFactor.factor_of_parts}
\end{lemma}

\begin{proof}
If $T$ is empty, $\Pi=\C$ with $p=p'$ the constant coefficient (and $P'=1$) is a factorisation. Otherwise fix $t_0\in T$ and let $K$
be the kernel of the linear map $P\mapsto I_{t_0}(P)$ (linear by the integrability of all kernels and linearity of the integral). By
the second hypothesis, $K$ is the kernel of $P\mapsto I_t(P)$ for every $t$. Put $\Pi=\shthetaPol(W,A)/K$, let $p$ be the quotient
map, $p'=p\circ\sigma$, and $J_t=\lambda(t)^{-1}\bar I_t$, where $\bar I_t:\Pi\to\{\text{functions}\}$ is induced by $I_t$; $\bar
I_t$, hence $J_t$, is injective. Then $I_t(P)=\lambda(t)J_t(p(P))$, and if the kernels of $P'$ are integrable,
$J_t(p'(P'))=\lambda(t)^{-1}I_t(\sigma P')=I'_t(P')$. Finally $p'(P'_0)\ne0$: otherwise $\sigma P'_0\in K$, so
$0=I_{t_0}(\sigma P'_0)=\lambda(t_0)I'_{t_0}(P'_0)$, contradicting $\lambda(t_0)\ne0\ne I'_{t_0}(P'_0)$.
\end{proof}

\begin{lemma}[Comparison of square constants with the top stage]\label{pf:sh-theta:compare}
Assume $F$ and Theorem~K. At large scale, for every real match $\varrho$ and every triple $X$ as in \ref{pf:sh-theta:c:grid} with a
factorisation with multiplier $\lambda$, there is $\kappa:T\to\C$ with $\kappa(t)\ne0$ and $\kappa(t)=\shthetaemb(\zeta_t)\kappa(\mu_0)$
for $Q(t)$-units $\zeta_t$, such that every square constant $c'$ at $(i,j,t)$ for $\varrho$ equals $\kappa(t)\lambda(t)$.
\formallabel{GlobalTheory.T.theta_square_norm_comparison}
\formallabel{GlobalTheory.SqCompare.theta_square_comparison}
\end{lemma}

\begin{proof}
Let $\varrho_0$ be the tied real match of \ref{pf:sh-theta:c:tied}, $l\ne0$ the number of \ref{pf:sh-theta:c:real} for
$(\varrho,\varrho_0)$ and $\kappa_0$ the function of \ref{pf:sh-theta:c:reg} for $\varrho_0$, and put $\kappa=l\kappa_0$. Then
$\kappa(t)=l\kappa_0(t)\ne0$ and $\kappa(t)=\shthetaemb(\zeta_t)\kappa(\mu_0)$. If $c'$ is a square constant at $(i,j,t)$, let $c_0$ be a
regularised constant for $\varrho_0$ at $(i,j,t)$ \ref{pf:sh-theta:c:tied}; then $c'=lc_0$ by \ref{pf:sh-theta:c:real} and
$c_0=\kappa_0(t)\lambda(t)$ by \ref{pf:sh-theta:c:reg}, so $c'=\kappa(t)\lambda(t)$.
\end{proof}

\begin{proof}[Proof of \cref{pf:sh-theta:five}]
Each statement is a combination of finitely many statements at large scale, so it suffices to verify the implications for
fixed data.

(T2) Let $X$ be a stage-$(p,q)$ triple. \Cref{pf:sh-theta:stageint} and the last statement of \ref{pf:sh-theta:c:stage} are the
hypotheses of \cref{pf:sh-theta:parts}, which gives the factorisation.

(T1) Let $X$ be the stage-$(p_j,q_j)$ triple of \ref{pf:sh-theta:c:grid} and $\lambda$ the multiplier of a factorisation
(T2); by \cref{pf:sh-theta:factor}(a) ($\mu_0\in T$) $c=\lambda$ is a constant of $X$, and \cref{pf:sh-theta:compare} gives
$\kappa$. The non-vanishing of the square is not used.

(T3) Let $Y$, $c$, $\phi_t$ be as in \ref{pf:sh-theta:c:seesaw}. The stage conditions for $X$ give
$\rank A=r_{i+1,j+1}-2(p_j-p)$ and $\rank A'=r_{ij}-2(q_j-q-1)$ with $q+1\le q_j$, hence
$2e=t+s-2(p_j-p)+2(q_j-q-1)\ge t+s-2p_j>0$, as $2p_j\le m<t$; so $e\ge1$, $e+1\ge2$, and
$L^{S_0}(\nu_t,e+1)=L_t(e+1)$ for all $t$ by \cref{pf:sh-theta:partialL}. Put $y(t)=\lambda(t)L_t(e+1)/c(t)$. Then
$\lambda(t)L_t(e+1)=y(t)c(t)$, and $y(t)\ne0$ because $\lambda(t)\ne0$, $c(t)\ne0$ and $L_t(e+1)=L^{S_0}(\nu_t,e+1)\ne0$ by
\ref{pf:sh-theta:c:seesaw}. By
\ref{pf:sh-theta:c:gauss}, $\phi_{\mu_0}=\shthetaemb(\gamma)\phi_t$ with $\gamma=b\prod_vg_v/g'_v$ as in
\cref{pf:sh-theta:units}. Substituting in the identity of \ref{pf:sh-theta:c:seesaw} and cancelling $\phi_t\ne0$ gives
$\lambda(t)L_t(e+1)c(\mu_0)=\shthetaemb(\gamma)\lambda(\mu_0)L_{\mu_0}(e+1)c(t)$; dividing by $c(t)c(\mu_0)$ gives
$y(t)=\shthetaemb(\gamma)y(\mu_0)$.

(T4) Let $Y$, $c$, $C$, $\zeta$, $\phi_t$ be as in \ref{pf:sh-theta:c:big}. The stage conditions give $e\ge2$: here
$2e=t+s-2(p_j-p-1)+2q_j\ge t+s-2p_j\ge4$, so $L^{S_0}(\nu_t,e)=L_t(e)$ for all
$t$ by \cref{pf:sh-theta:partialL}. Since $\lambda(t)\ne0$, the identity of \ref{pf:sh-theta:c:big} shows $C\ne0$,
$\phi_t\ne0$ and $L_t(e)c(t)\ne0$ for every $t$. Put $y(t)=\lambda(t)/(L_t(e)c(t))$, so $y(t)\ne0$ and
$\lambda(t)=y(t)L_t(e)c(t)$. By \ref{pf:sh-theta:c:big}, $y(t)=C\shthetaemb(\zeta)\phi_t$ and $y(\mu_0)=C\phi_{\mu_0}$; by
\ref{pf:sh-theta:c:gauss} $\phi_{\mu_0}$ is a step multiple of $\phi_t$, and $\phi_{\mu_0}\ne0$. By
\cref{pf:sh-theta:unitlemma}(c), $y(t)=C\shthetaemb(\zeta)\phi_t$ is a step multiple of $C\phi_{\mu_0}=y(\mu_0)$.

(T5) Apply the four statements of \ref{pf:sh-theta:c:bottom} in turn: $Y'$ on the spaces of $X$ has the constant $c$; then
$Z_1$ with the lattices of $X$ has the constant $v_1c$; then $Z_2$ with the lattices and frames of $X$ has the constant
$v_2v_1c$; finally $X$ has the constant $\kappa v_2v_1c$ for some $\kappa\ne0$. Put $v=\kappa v_2v_1$. Then $v(t)\ne0$ and
$v(t)=\shthetaemb(\zeta_{2,t}\zeta_{1,t})v(\mu_0)$, and $\zeta_{2,t}\zeta_{1,t}$ is a $Q(t)$-unit by
\cref{pf:sh-theta:unitlemma}(a). The hypotheses $F$ and Theorem~K are not used for (T5).
\end{proof}

\subsubsection{The descent}

\begin{lemma}[Descent along the tower]\label{pf:sh-theta:descent}
Fix data at which \ref{pf:sh-theta:T2}, \ref{pf:sh-theta:T3} and \ref{pf:sh-theta:T4} hold, a square $(i,j)$ and
$L$-functions $L_t$ of $\nu_H(t,i,j)$ for all $t$. Let $X$ be a stage-$(p,q)$ triple with a constant $c$ and
$\rank A=\rank A'+2e$ with $p\le e$. Then there are a stage-$(0,0)$ triple $Y$ with a constant $c_Y$, functions
$u_k:T\to\C$ ($k\in\N$) with $u_k(t)\ne0$ and $u_k(t)$ a step multiple of $u_k(\mu_0)$, and $\kappa\ne0$, such that for all
$t$
\begin{equation}\label{eq:pf:sh-theta:descent}
c(t)\prod_{k<q}L_t(e+k+1)=\kappa\Bigl(\prod_{k<p+q}u_k(t)\Bigr)c_Y(t)\prod_{k<p}L_t(e+q-k).
\end{equation}
\formallabel{GlobalTheory.SqTower.descent}
\end{lemma}

\begin{proof}
Induction on $q$, and for $q=0$ induction on $p$.

\emph{$q=0$, $p=0$.} Take $Y=X$, $c_Y=c$, $u_k=1$ and $\kappa=1$; both sides of \eqref{eq:pf:sh-theta:descent} equal $c(t)$,
and $1$ is a step multiple of itself.

\emph{$q=0$, $p\to p+1$.} Let $X$ be a stage-$(p+1,0)$ triple with constant $c$, $\rank A=\rank A'+2e$ and $p+1\le e$. By
(T2) $X$ has a factorisation with some multiplier $\lambda$, and $c=\kappa_1\lambda$ with $\kappa_1\ne0$ by
\cref{pf:sh-theta:factor}(b). By (T4) there are a stage-$(p,0)$ triple $Y_1$ with constant $c_1$ and $y$ with
$y(t)\ne0$, $y(t)$ a step multiple of $y(\mu_0)$, and $\lambda(t)=y(t)L_t(e)c_1(t)$. Comparing the stage conditions of $X$ and
$Y_1$ (the ranks of $A$ and $A'$ are $r_{i+1,j+1}-2(p_j-p-1)$, $r_{ij}-2q_j$ for $X$ and $r_{i+1,j+1}-2(p_j-p)$,
$r_{ij}-2q_j$ for $Y_1$), the spaces of $Y_1$ satisfy $\rank B=\rank B'+2(e-1)$, and $p\le e-1$. The induction hypothesis
for $(Y_1,c_1,e-1)$ gives $Y$, $c_Y$, $u$ and $\kappa$ with
$c_1(t)=\kappa\prod_{k<p}u_k(t)\,c_Y(t)\prod_{k<p}L_t(e-1-k)$. Put $u'_0=y$, $u'_{k+1}=u_k$ and $\kappa'=\kappa_1\kappa$.
Then
\begin{align*}
c(t)&=\kappa_1y(t)L_t(e)c_1(t)=\kappa'\prod_{k<p+1}u'_k(t)\,c_Y(t)\,L_t(e)\prod_{k<p}L_t(e-(k+1))\\
&=\kappa'\prod_{k<p+1}u'_k(t)\,c_Y(t)\prod_{k<p+1}L_t(e-k),
\end{align*}
which is \eqref{eq:pf:sh-theta:descent} for $(p+1,0)$.

\emph{$q\to q+1$.} Let $X$ be a stage-$(p,q+1)$ triple with constant $c$, $\rank A=\rank A'+2e$, $p\le e$. As before
$c=\kappa_1\lambda$ for the multiplier $\lambda$ of a factorisation (T2). By (T3) there are a stage-$(p,q)$ triple $Y_1$
with constant $c_1$ and $y$ with $y(t)\ne0$, $y(t)$ a step multiple of $y(\mu_0)$ and $\lambda(t)L_t(e+1)=y(t)c_1(t)$. Now
$\rank A'$ drops by $2$ and $\rank A$ is unchanged, so $\rank B=\rank B'+2(e+1)$, and $p\le e+1$. The induction hypothesis
for $(Y_1,c_1,e+1)$ gives
$c_1(t)\prod_{k<q}L_t(e+k+2)=\kappa\prod_{k<p+q}u_k(t)\,c_Y(t)\prod_{k<p}L_t(e+1+q-k)$. With $u'$ and $\kappa'$ as before,
\begin{align*}
c(t)\prod_{k<q+1}L_t(e+k+1)&=\kappa_1\lambda(t)L_t(e+1)\prod_{k<q}L_t(e+k+2)=\kappa_1y(t)c_1(t)\prod_{k<q}L_t(e+k+2)\\
&=\kappa'\prod_{k<p+q+1}u'_k(t)\,c_Y(t)\prod_{k<p}L_t(e+(q+1)-k),
\end{align*}
which is \eqref{eq:pf:sh-theta:descent} for $(p,q+1)$.
\end{proof}

\subsubsection{The tower formula and bottom triples}

Recall (\cref{pf:fit:S:LNe}) that $\sfname{LNe}$ asserts at large scale: for every square $(i,j)$ and type $t$ there is an
$L$-function $\Lambda$ of $\nu_H(t,i,j)$ with $\Lambda(s)\ne0$ for all integers $1-e_0\le s\le h$, and for every
$c\in\{\mathrm{none},\mathrm{false},\mathrm{true}\}$ an $L$-function of the character $\mathrm{fam}(G,\mu_0,i,j,c)$ which does
not vanish at the integers $s\ge1$.

\begin{theorem}[The Witt-tower formula; \ref{pf:fit:S:Tower}]\label{pf:sh-theta:towerthm}
Assume $F$, Theorem~K and $\sfname{LNe}$. Then $\sfname{Tower}$ holds: at large scale, for every Weil datum $W$ with standard
character, grid choice $S$ satisfying $\sfname{NumbersIn}$ at $\mu_0$, real match $\varrho$, square $(i,j)$ and type $t$,
every bottom triple $X$ with a constant $c_0$, all square constants $c$ at $(i,j,t)$ and $c_\circ$ at $(i,j,\mu_0)$ of
non-zero squares, and all $L$-functions $\Lambda$ of $\nu_H(t,i,j)$ and $\Lambda_\circ$ of $\nu_H(\mu_0,i,j)$, there are
$u_1,\dots,u_{p_j+q_j}\in\C$, each a step multiple of $1$ at $t$, and a $Q(t)$-unit $\zeta$ such that
\begin{equation}\label{eq:pf:sh-theta:tower}
\frac c{c_\circ}\prod_{k<q_j}\frac{\Lambda(e_T+k+1)}{\Lambda_\circ(e_T+k+1)}
=\shthetaemb(\zeta)\Bigl(\prod_ku_k\Bigr)\frac{c_0(t)}{c_0(\mu_0)}\prod_{k<p_j}\frac{\Lambda(e_T+q_j-k)}{\Lambda_\circ(e_T+q_j-k)}.
\end{equation}
\formallabel{GlobalTheory.SqTower.tower_alt_bis}
\end{theorem}

\begin{proof}
At large scale the five statements of \cref{pf:sh-theta:five} and $\sfname{LNe}$ hold simultaneously; fix such data and
$W,S,\varrho,(i,j),t,X,c_0,c,c_\circ,\Lambda,\Lambda_\circ$ as in the statement.

\emph{$L$-values.} By $\sfname{LNe}$ choose for every type $t'$ an $L$-function $L_{t'}$ of $\nu_H(t',i,j)$ which does not vanish
at the integers in $[1-e_0,h]$. By uniqueness, $\Lambda=L_t$ and $\Lambda_\circ=L_{\mu_0}$. Since $p_j\le e_T$ and
$e_T+q_j\le h$, all arguments $e_T+k+1$ ($k<q_j$) and $e_T+q_j-k$ ($k<p_j$) lie in $[1,h]$, so the two products of values of
$L_{\mu_0}$ in \eqref{eq:pf:sh-theta:tower} are non-zero.

\emph{Top stage.} By (T1) there are a stage-$(p_j,q_j)$ triple $X_1$ with a constant $c_1$ and $\kappa_T$ with
$\kappa_T(t')\ne0$, $\kappa_T(t')=\shthetaemb(\zeta_{T,t'})\kappa_T(\mu_0)$, such that $c=\kappa_T(t)c_1(t)$ and
$c_\circ=\kappa_T(\mu_0)c_1(\mu_0)$. By (T2) and \cref{pf:sh-theta:factor}(b), $c_1=\kappa'\lambda_1$ with $\kappa'\ne0$ and
$\lambda_1$ the multiplier of a factorisation, so $c_1(t')\ne0$ for all $t'$. The stage conditions give $\rank A_1=r_{i+1,j+1}$
and $\rank A_1'=r_{ij}$, i.e.\ $\rank A_1=\rank A'_1+2e_T$.

\emph{Descent.} \Cref{pf:sh-theta:descent} for $(X_1,c_1)$ with $(p,q,e)=(p_j,q_j,e_T)$ (note $p_j\le e_T$) gives a
stage-$(0,0)$ triple $Y$ with a constant $c_Y$, $u_k$ and $\kappa\ne0$ with
\begin{equation}\label{eq:pf:sh-theta:top}
c_1(t')\prod_{k<q_j}L_{t'}(e_T+k+1)=\kappa\prod_{k<p_j+q_j}u_k(t')\;c_Y(t')\prod_{k<p_j}L_{t'}(e_T+q_j-k)\qquad(t'\in T).
\end{equation}

\emph{Bottom.} By (T5) for $X$ and $Y$ there is $v$ such that $vc_Y$ is a constant of $X$, with $v(t')\ne0$ and
$v(t)=\shthetaemb(\zeta_v)v(\mu_0)$ for a $Q(t)$-unit $\zeta_v$. By (T2) the bottom triple $X$, which is a stage-$(0,0)$ triple,
has a factorisation with some multiplier $\lambda_0$, and by \cref{pf:sh-theta:factor}(b) $c_0=a\lambda_0$ and
$vc_Y=b_0\lambda_0$ with $a,b_0\ne0$; hence $c_0(t')=(a/b_0)\,v(t')c_Y(t')$ for all $t'$.

\emph{Non-vanishing at $\mu_0$.} The left side of \eqref{eq:pf:sh-theta:top} at $t'=\mu_0$ is non-zero, so $c_Y(\mu_0)\ne0$;
also $u_k(\mu_0)\ne0$ and $v(\mu_0)\ne0$, so $c_0(\mu_0)\ne0$.

\emph{Conclusion.} Put $u'_k=u_k(t)/u_k(\mu_0)$, a step multiple of $1$ by \cref{pf:sh-theta:unitlemma}(b), and
$\zeta=\zeta_{T,t}\zeta_v^{-1}$, a $Q(t)$-unit by \cref{pf:sh-theta:unitlemma}(a) ($\zeta_v\ne0$ since $v(t)\ne0$). Using
$c/c_\circ=\kappa_T(t)c_1(t)/(\kappa_T(\mu_0)c_1(\mu_0))=\shthetaemb(\zeta_{T,t})\,c_1(t)/c_1(\mu_0)$ and
\eqref{eq:pf:sh-theta:top} at $t$ and $\mu_0$,
\[
\frac c{c_\circ}\prod_{k<q_j}\frac{L_t(e_T+k+1)}{L_{\mu_0}(e_T+k+1)}
=\shthetaemb(\zeta_{T,t})\prod_k\frac{u_k(t)}{u_k(\mu_0)}\cdot\frac{c_Y(t)}{c_Y(\mu_0)}
\prod_{k<p_j}\frac{L_t(e_T+q_j-k)}{L_{\mu_0}(e_T+q_j-k)},
\]
and $c_Y(t)/c_Y(\mu_0)=\shthetaemb(\zeta_v)^{-1}c_0(t)/c_0(\mu_0)$ by the bottom step. This is \eqref{eq:pf:sh-theta:tower}.
\end{proof}

\begin{remark}\label{pf:sh-theta:towerrem}
Of $\sfname{LNe}$ only two consequences are used: the existence of $L$-functions of $\nu_H(t',i,j)$ for all types $t'$,
and the non-vanishing of the $L$-function of $\nu_H(\mu_0,i,j)$ at the integers $1\le s\le h$ at the base type. The
non-vanishing of the squares enters only through (T1), whose proof does not use it.
\end{remark}

The second statement, \ref{pf:fit:S:TriEx}, asserts the existence of bottom triples with constants.

\begin{theorem}[Bottom triples with constants; \ref{pf:fit:S:TriEx}]\label{pf:sh-theta:triex}
Assume Howe duality at the real place ($\sfname{HoweDualityReal}$), $F$, Theorem~K and $\sfname{LNe}$. Then $\sfname{TriEx}$
holds: at large scale, for every Weil datum $W$ with standard character and every square $(i,j)$ there are a bottom triple
$X$ and a constant $c_0$ of $X$.
\formallabel{GlobalTheory.SqTop.triEx_holds_alt}
\end{theorem}

\begin{proof}
We use three results proved elsewhere. Bottom triples exist for every Weil datum and square
(\lref{GlobalTheory.RowTest.exists_bottomTri}: the standard riggings, Weil representations with the required
splitting characters, vacuum vectors and the theta functions of the unit-character sections are constructed explicitly), and
their kernels are integrable, the spaces being definite (\lref{GlobalTheory.RowTest.IsBottomTri.integrable}). Under the
hypotheses, at large scale, for every bottom triple $X$ there are a vector space $\Pi$, a surjective linear map $p$, a linear
map $p'$ with $p'(P'_1)\ne0$ for some $P'_1$, and linear maps $J_t$ each of which is zero or injective, such that for every $t$
there is $\lambda_t$ with $I_t(P)=\lambda_tJ_t(p(P))$ for all $P$, and $I'_t(P')=J_t(p'(P'))$ for all $P'$
(\lref{GlobalTheory.T.bottom_theta_norm_eq_prod_local_factors}); and for every $t$ some $I_t(P)$ and some $I'_t(P')$ are
non-zero functions (\lref{GlobalTheory.T.first_occurrence_theta_lift_ne_zero}: the bottom theta lifts are non-zero).

Let $X$ be a bottom triple. For every $t$, $J_t$ is injective: if $J_t=0$, then $I'_t(P')=0$ for all $P'$, which is
excluded. And $\lambda_t\ne0$: otherwise $I_t(P)=0$ for all $P$. All kernels are integrable, so $(\Pi,p,p',J)$ is a
factorisation of $X$ with multiplier $t\mapsto\lambda_t$ (with the polynomial $P'_1$), and $T\ni\mu_0$ is non-empty; by
\cref{pf:sh-theta:factor}(a), $c_0=\lambda$ is a constant of $X$. The thresholds are the maxima of those of the two cited
statements.
\end{proof}

% ---------------------------------------------------------------------------
\subsection{Doubling integrals, zeta integrals and level indices}\label{sec:pf:sh-theta:doubling}

This subsection proves the inputs on doubling integrals used in \cref{pf:liftinj:euler} (the unramified doubling
integral at inert places), \cref{pf:bundles:thm:chain} and \cref{pf:fs:stepUp} (the Rallis inner product formula), and in
\cref{pf:onto3:adelc} (the place models, the finite zeta integral of a theta lift and the product of the level index with
the norm of the level vector).

\subsubsection{The unramified doubling integral at an inert place}

Let $L$ be a complete discretely valued field with uniformiser $\varpi$ and $\sigma$ an involution of $L$ such that
$L/L^\sigma$ is unramified quadratic: $\sigma$ preserves the valuation, $\sigma\varpi=\varpi$, some integral $x$ has
$\sigma x-x$ a unit, and the residue field of $L$ has $q^2$ elements with $q\ge2$. Let $G_n=\Uni(n)$ be the unitary group of
the antidiagonal hermitian form on $L^n$, $K_n$ its subgroup of integral points, $G_n^\square$ the unitary group of the doubled
form on $L^{2n}$, with Siegel parabolic subgroup $P^\square=M^\square N^\square$ ($a(p)\in L^\times$ the determinant of the
Levi component) and maximal compact subgroup $K^\square$, and $\iota(g)=\iota(g,1)$ the doubling embedding of $G_n$. For an
unramified character $\chi$ of $L^\times$ and $s\in\C$, a \emph{spherical section} is a function $\Phi$ on $G^\square_n$ with
$\Phi(pk)=\chi(a(p))\,q^{-2\ord(a(p))(s+n/2)}$ for $p\in P^\square$, $k\in K^\square$. For $\alpha\in(\C^\times)^{\lfloor n/2\rfloor}$
a \emph{zonal spherical function} with parameters $\alpha$ is a bi-$K_n$-invariant function $w$ with $w(1)=1$ on which the
spherical Hecke algebra acts through the Satake transform at $\alpha$. Put
\[
X=\chi(\varpi)q^{-(2s+n)},\qquad y=\chi(\varpi)q^{-(2s+1)},
\]
\[
Z_n(X,y,\alpha)=\prod_{j=0}^{n-1}\bigl(1-(-q)^jX\bigr)\prod_{i}\frac1{(1-y\alpha_i)(1-y\alpha_i^{-1})}\cdot(1-y)^{-[n\text{ odd}]}.
\]

\begin{theorem}[Unramified doubling integral at an inert place]\label{pf:sh-theta:inert}
In this situation let $\mu$ be the Haar measure of $G_n$ with $\mu(K_n)=1$.
\begin{enumerate}[label=(\alph*)]
\item If $|y|<1$ and $|y\alpha_i^{\pm1}|<1$ for all $i$, then $g\mapsto\Phi(\iota(g))w(g)$ is integrable and
\[
\int_{G_n}\Phi(\iota(g))\,w(g)\,d\mu(g)=Z_n(X,y,\alpha);
\]
if moreover $\chi(\varpi)=(-1)^m$, this equals $L(s+\frac12,\mathrm{BC}(\alpha)\otimes\chi)/d_{n,m}(s)$, where
$L(z,\mathrm{BC}(\alpha)\otimes\chi)=\prod_i(1-\chi(\varpi)\alpha_iq^{-2z})^{-1}(1-\chi(\varpi)\alpha_i^{-1}q^{-2z})^{-1}
(1-\chi(\varpi)q^{-2z})^{-[n\text{ odd}]}$ is the Euler factor of the base change parameters and
$d_{n,m}(s)=\prod_{i=0}^{n-1}\bigl(1-(-1)^{m+i}q^{-(2s+n-i)}\bigr)^{-1}$.
\item If $n\le1$, the integral converges and equals $1$.
\end{enumerate}
\formallabel{GlobalTheory.InertZeta.unramDoublingInert}
\end{theorem}

The restatement \lref{GlobalTheory.T.unramifiedDoublingIntegral_inert_place} of this theorem, with the same hypotheses and
conclusion, is the form in which it is applied in \cref{sec:pf:stageU}.

This is the unramified computation of the doubling zeta integral of Piatetski-Shapiro and Rallis
\cite{PiatetskiShapiroRallis1987,GelbartPiatetskiShapiroRallis1987} at an inert place (see also \cite{LapidRallis2005} for the
normalising factors); it is carried out in the formal development by an explicit count of lattices.

\begin{proof}
\emph{The section on $\iota(G_n)$.} Every $g\in G_n$ has an elementary divisor \emph{length} $\ell(g)\in\N$: write
$(U+Vg)C=1$ with $U,V,C,gC$ integral and let $\ell(g)$ be the least $k$ with $v(\det C)=-k$ for such a fraction
(\lref{GlobalTheory.T.matrixLengthExists_of_elementaryDivisors}: a fraction realising $\ell(g)$ exists, by elementary
divisors). For such a fraction, the doubled lattice has a primitive isotropic frame completing $(Cx,gCx)$
(\lref{GlobalTheory.T.doubledLattice_primitive_isotropic_frame_completion}: there is $k\in K^\square$ with
$k(x,x)=(Cx,gCx)$ for all $x$), and evaluating the section with this decomposition gives $\Phi(\iota(g))=X^{\ell(g)}$
(\lref{GlobalTheory.InertZeta.section_eq}). The length is
bi-$K_n$-invariant (\lref{GlobalTheory.InertZeta.len_mat_mul}) and vanishes on $K_n$
(\lref{GlobalTheory.InertZeta.len_of_mem}).

\emph{Reduction to a coset sum.} The integrand $X^{\ell(g)}w(g)$ is right $K_n$-invariant and $K_n$ is open compact of volume
$1$; hence the integral exists and equals $Z$ as soon as $\sum_{x\in G_n/K_n}X^{\ell(x)}w(x)$ converges to $Z$
(\lref{GlobalTheory.T.integral_right_invariant_eq_sum_open_subgroup_cosets}).

(b) If $n\le1$, $G_n=K_n$ (\lref{GlobalTheory.T.unitary_rank_le_one_integral_points}), so the sum has the single term
$X^0w(1)=1$.

(a) By the Satake theory for unramified unitary groups there is a spherical vector $\varphi$ of the unramified principal series
with parameters $\alpha$, normalised by $\varphi(1)=1$ (\lref{Formalism.UG.Final.SphericalExistsU}). The \emph{generating
series} of the length,
\[
\sum_{x\in G_n/K_n}X^{\ell(x)}\varphi(x)=Z_n(X,y,\alpha)\qquad(\text{with }yq=Xq^n,\ |y|<1,\ |y\alpha_i^{\pm1}|<1),
\]
is \lref{GlobalTheory.InertZeta.hasSum_spherical}; it is proved by induction on $\lfloor n/2\rfloor$ through the Iwasawa
decomposition with respect to the parabolic subgroup stabilising an isotropic line, from eleven local statements proved in the
formal development for this purpose: existence of fractions, the lengths of elements $mu$ with $m$ in the Levi subgroup and $u$
in the unipotent radical in terms of a smallness condition, a count of the cosets of the radical satisfying it, a corner lemma,
a generating function for counts of the shape $c\,q^{2N\lfloor j/2\rfloor+j}$, surjectivity and disjointness of the Iwasawa cells,
existence of Levi elements, transitivity of spherical induction, $G_n=K_n$ for $n\le1$, and the geometric summation over the
Levi variable (\lref{GlobalTheory.T.leviRadicalProductLength_of_lattice_steps}, \lref{GlobalTheory.T.smallRadicalCosetCount_of_heisenberg}, \lref{GlobalTheory.T.small_corner_of_heisenberg_duality}, \lref{GlobalTheory.T.length_count_generating_series_formula}, \lref{GlobalTheory.T.unitaryIwasawaCosetSurj_of_leviElementExists}, \lref{GlobalTheory.T.unitary_iwasawa_cells_disjoint}, \lref{GlobalTheory.T.exists_unitary_block_diagonal_levi_element}, \lref{GlobalTheory.T.sphericalInductionTransitivity_of_leviExists}, and \lref{GlobalTheory.T.unitary_levi_summation_geometric_series}).
The functions $\varphi$ and $w$ are both bi-$K_n$-invariant (for $\varphi$, after averaging over $K_n$) and have the same Hecke
eigenvalues (\lref{Formalism.UG.Final.SatakeInert}); since $X^{\ell}$ is bi-$K_n$-invariant, the coset sums of $X^\ell\varphi$
and of $X^\ell w$ agree (\lref{GlobalTheory.T.biInvariant_replace_same_double_coset_sums}). With $yq=Xq^n$ (which holds for
the values above) this gives the first formula. For the second, if $\chi(\varpi)=(-1)^m$ then
$1-(-q)^jX=1-(-1)^{m+j}q^{-(2s+n-j)}$, so the first product is $d_{n,m}(s)^{-1}$, and the remaining factors are the Euler factor
at $z=s+\frac12$, since $\chi(\varpi)q^{-2z}=y$.
\end{proof}

\subsubsection{The Rallis inner product formula}

We use the doubling pairs, the integrand $I_s$ and the unramified Euler product (ER) of \cref{pf:liftinj:ERdef}, and the Rallis
inner product formula of \cref{pf:bundles:def:rallis}, which is the statement $\sfname{HeckeForm}(D)$. Write
$s_0=(\dim B-\dim A)/2=t/2$ and, for a doubling pair and cusp forms $u_1,u_2$ on $\Uni(A)$, $Z(s)=\int I_s(g)\,dg$ and
$Z_{S}(s)=\int_{\mathcal O_S}I_s(g)\,dg$, the integral over the elements integral outside $S$. The \emph{global form}
$\sfname{GlobalForm}(D)$ asserts: for every self-dual lattice $L$ of $A$, every doubling pair with $B$ of Witt index $0$ or
$<t$, Haar measures and Fock test vectors $\Phi_1,\Phi_2$ and $u_1,u_2\in\mathcal A_0(A)$, the theta lifts $\theta_{\Phi_i}(u_i)$
converge and are square integrable, $I_s$ is integrable for $\Re s>\dim A/2$, and $Z$ has a holomorphic value $z$ at $s_0$ with
$\mathrm{pet}(\theta_{\Phi_1}(u_1),\theta_{\Phi_2}(u_2))=\vol(B)\,z$ (the Siegel--Weil formula and the doubling identity).

\begin{lemma}[From the global form and the Euler product to the inner product formula]\label{pf:sh-theta:ipf}
For every global datum $D$, $\sfname{GlobalForm}(D)$ and (ER) for $D$ imply $\sfname{HeckeForm}(D)$.
\formallabel{GlobalTheory.T.heckeInnerProductFormula_of_doublingEuler}
\formallabel{GlobalTheory.Core.Rk.IP.heckeForm_of_global_euler}
\end{lemma}

\begin{proof}
Fix the data of \cref{pf:bundles:def:rallis}: in particular $s_1<s_0$ such that $I_{s_1}$ is integrable on $\mathcal O_S$. By
$\sfname{GlobalForm}$ the lifts converge and there is a holomorphic value $z$ of $Z$ at $s_0$ with $\mathrm{pet}=\vol(B)z$. On
$\mathcal O_S$ the height $\mathrm{ht}_L$ is continuous with values in $(0,1]$ (\lref{GlobalTheory.T.doubledLatticeHeightFun_properties}),
and $I_s=c_\omega r\,\mathrm{ht}_L^{s-s_0}$ is a \emph{flat family}: since $|\mathrm{ht}_L^{s-s_0}|\le\mathrm{ht}_L^{s_1-s_0}$ for
$\Re s\ge s_1$, the integrals $Z_S(s)$ converge absolutely for $\Re s\ge s_1$ and define a holomorphic function on $\Re s>s_1$
(\lref{GlobalTheory.T.flatFamilyZeta_convergence_holomorphy}). In particular, as $s_1<s_0$, $c_\omega r=I_{s_0}$ is integrable on
$\mathcal O_S$. By (ER), for $\Re s>\dim A+1$ (where $I_s$ is integrable by $\sfname{GlobalForm}$),
\[
Z(s)=\frac{L^S(s+\frac12,\mathrm{par}\times\chi_B)}{d^S(s)}\,Z_S(s).
\]
The function $d^S$ is analytic and non-vanishing on $\Re s>0$ (\lref{GlobalTheory.T.partialNormalizingFactorProduct_analytic_ne_zero}),
$Z_S$ is holomorphic on $\Re s>s_1$, and $l$ is a holomorphic value of $L^S(s+\frac12)$ at $s_0$. Since both sides agree on a
half-plane, their holomorphic values at $s_0$ agree: $z=\frac{l}{d^S(s_0)}Z_S(s_0)$
(\lref{GlobalTheory.T.holomorphicContinuationValue_eq_div_mul}, the identity theorem). Hence
$\mathrm{pet}=\vol(B)\frac{l}{d^S(t/2)}\int_{\mathcal O_S}c_\omega(g)r(g)\,dg$, which is the formula of \cref{pf:bundles:def:rallis}.
\end{proof}

\begin{theorem}[Rallis inner product formula]\label{pf:sh-theta:heckeform}
The existence of Weil representations for unitary dual pairs, rapid decay of cusp forms, the Weil--Siegel formula, the
Siegel--Weil formula in the convergent range, the local zeta integrals of the doubling method and the unramified doubling
computation for $\GL_n$ ($\sfname{UnramifiedDoublingGL}$) imply $\sfname{HeckeForm}(D)$ for every global datum $D$.
\formallabel{GlobalTheory.Core.Rk.IP.heckeForm_of_doubling'}
\end{theorem}

\begin{proof}
\emph{Global form.} The first five hypotheses imply $\sfname{GlobalForm}(D)$
(\lref{GlobalTheory.T.globalForm_of_siegel_weil}): the Siegel--Weil formula for the Fock doubled pair (from the Weil--Siegel
formula and \cite{Ichino2004}) gives $\sfname{GlobalForm}$ at the frame lattice
(\lref{GlobalTheory.T.globalFormAt_frameLat_of_siegelWeil}), and it is transferred to every self-dual lattice
(\lref{GlobalTheory.T.global_doubling_zeta_form_selfDual_lattice_transfer}).

\emph{Euler product.} At the split places the local unramified zeta integrals are computed from $\sfname{UnramifiedDoublingGL}$
(\lref{GlobalTheory.T.unramifiedDoublingIntegral_splitPlace_adelic}) and the spherical parameters are bounded
(\lref{GlobalTheory.T.sphericalParameterBound_split_place}); at the non-split places both follow from
\cref{pf:sh-theta:inert} (\lref{GlobalTheory.T.sphericalParameterBound_nonsplit_places} and \lref{GlobalTheory.T.unramified_doubling_integral_inert_adelic}).
Combining the two kinds of places gives the local statements at all places outside $S$
(\lref{GlobalTheory.Core.Rk.IP.paramBoundOn_all_of_split_inert} and \lref{GlobalTheory.Core.Rk.IP.localZetaOn_all_of_split_inert}). With
the spherical factorisation of the doubling integrand, which follows from the existence of Weil representations
(\lref{GlobalTheory.T.doublingSphericalFactorization_of_weilDualPair}), and rapid decay, these give (ER) for $D$
(\lref{GlobalTheory.T.doubling_eulerProduct_farRight}). \Cref{pf:sh-theta:ipf} concludes.
\end{proof}

\subsubsection{Place models}

Let $D$ be a datum. For hermitian spaces $S,T$, a Weil representation $\omega$ of $(S,T)$, Cartan data at the real places,
self-dual lattices $\Lambda_S,\Lambda_T$ and a \emph{base vector} $f_0$ (the theta function which is the lattice vacuum at the
finite places and the Fock vacuum at the real places), a family $x=(x_v)$ indexed by the places is \emph{restricted} if each $x_v$
is a local theta vector at $v$ and $x_v=f_0$ for almost all $v$, and a theta function $f$ is the \emph{pure tensor} of $x$ if its
inner products with all translates are the products of the local ones. The inner product of theta functions is the $L^2$-product
on a fundamental domain of covolume one for the rational points.

\begin{definition}[Place models]\label{pf:sh-theta:pm}
$\sfname{PlaceModels}(D)$ is the conjunction of the following twelve statements, for all data as above:
\begin{enumerate}[label=(P\arabic*),leftmargin=*]
\item\label{pf:sh-theta:P1} (unitarity) theta functions form a vector space, the inner product is defined, non-degenerate,
independent of the covolume-one measure, preserved by Heisenberg translations and by $\omega$, and $p\mapsto\|\omega(p)f_1-f_2\|$
is continuous;
\item\label{pf:sh-theta:P2} (restricted tensor product) local theta vectors are preserved by the local Heisenberg group and the
local unitary groups; every restricted family has a unique pure tensor; the span of pure tensors is dense; Heisenberg
translations and $\omega$ act on pure tensors place by place;
\item\label{pf:sh-theta:P3} (matrix coefficients) for pure tensors $f_1,f_2$ of $x_1,x_2$, the matrix coefficient
$\langle\omega(g,h)f_1,f_2\rangle$ is the product over all places $v$ of the local matrix coefficients
$\langle\omega_v(g_v,h_v)x_{1,v},x_{2,v}\rangle$, finitely many of which are different from $1$;
\item\label{pf:sh-theta:P4} (split places) at a split finite place $v$, a Schr\"odinger model for the first half of the
symplectic space: a bijection from the space of Schwartz--Bruhat functions on the matrix space over $F_v$ onto the space of
local theta vectors, which multiplies inner products by the volume of the image of the half lattice, and on which the Levi
group acts linearly up to the splitting characters;
\item\label{pf:sh-theta:P5} the same for the second (adjoint) half;
\item\label{pf:sh-theta:P6} (split lattices) at split places the images of the half lattices are the boxes
$\{y:\ v(y_{ij})\le v(a)\}$ for an extremal element $a$ of the host;
\item\label{pf:sh-theta:P7} (vacua) at the non-split places where the characters are unramified, the stabilisers of the lattices
fix the lattice vacuum;
\item\label{pf:sh-theta:P8} the same at the split places;
\item\label{pf:sh-theta:P9} (Fock model) at every real place, for Fock vacua $f$ and polynomials $P,Q$: the Fock pure tensors are
theta functions, $\langle\mathrm{fock}(f,P),\mathrm{fock}(f,Q)\rangle=\langle P,Q\rangle_{\mathrm{Fock}}\|f\|^2$, the maximal compact
subgroups act by determinant characters composed with a linear substitution of the variables, and the non-compact part of the Lie
algebras acts by second order operators (creation and annihilation pairs);
\item\label{pf:sh-theta:P10} (Fock vectors are pure tensors) $\mathrm{fock}(f,\prod_\tau P_\tau)$ is the pure tensor of the family
obtained from that of $f$ by replacing the archimedean components with the local Fock vectors;
\item\label{pf:sh-theta:P11} (definite theta lifts) the theta vectors of the definite theta lifts are pure tensors of restricted
families which equal the base vector outside the level and twist places, and whose components at these places are given, in the
Schr\"odinger model of the adjoint half, by local newform functions: of principal level at the level places, of full level at
the twist places;
\item\label{pf:sh-theta:P12} (chain steps) similarly, the step vectors of theta chains are pure tensors of restricted families
which equal the base vector outside the level and twist places, with components given, in the Schr\"odinger model of the first
half, by the local newform functions of the local levels of the step at the level places and by an explicit function at the twist
places.
\end{enumerate}
\formal{GlobalTheory.IndefiniteRowSideFullHecke.RestrictedTensorSplitting.PlaceModels}
\end{definition}

\begin{theorem}[Place models]\label{pf:sh-theta:placemodels}
$\sfname{PlaceModels}(D)$ holds for every global datum $D$.
\formallabel{GlobalTheory.PlaceTensor.placeModels}
\end{theorem}

\begin{proof}
We indicate the source of each component. \ref{pf:sh-theta:P1}: from the existence of a covolume-one measure, the hermitian
properties of the theta inner product, the Heisenberg laws, the unitarity of $\omega$ and the continuity of
$p\mapsto\|\omega(p)f_1-f_2\|$ (\lref{GlobalTheory.PlaceTensor.unitary}).
\ref{pf:sh-theta:P2}: \lref{GlobalTheory.PlaceTensor.restrictedTensor}, from the stability of local factors, the existence and
uniqueness of pure tensors of restricted families, the placewise action, and the density of the span of Heisenberg translates of
the base vector (\lref{GlobalTheory.T.heisenberg_translates_span_dense_thetaFunctions}).
\ref{pf:sh-theta:P3}: the placewise action and the product formula for inner products of pure tensors
(\lref{GlobalTheory.PlaceTensor.coeffProduct}). \ref{pf:sh-theta:P4} and \ref{pf:sh-theta:P5}: the matrix coordinates on the
two halves at a split place and the Schr\"odinger model
(\lref{GlobalTheory.PlaceTensor.splitModel} and \lref{GlobalTheory.PlaceTensor.splitModelAdj}).
\ref{pf:sh-theta:P6}: \lref{GlobalTheory.T.splitPlace_thetaLattice_halves_selfDual_frames}. \ref{pf:sh-theta:P7} and
\ref{pf:sh-theta:P8}:
\lref{GlobalTheory.T.inert_place_stabilizers_fix_weil_lattice_vacuum} and \lref{GlobalTheory.T.splitPlace_latticeStabilisers_fix_weilVacuum}.
\ref{pf:sh-theta:P9}: \lref{GlobalTheory.PlaceTensor.fockModel}: the character of the maximal compact subgroup on the Fock
vacuum, the Fock inner product, the action of the maximal compact subgroup by substitution and of the non-compact part by
quadratic operators, the last because the matrix coefficients of the vacuum are even
(\lref{GlobalTheory.T.fockVacuum_matrixCoefficient_even}; a function with an even matrix coefficient has derivative $0$ at
$0$). \ref{pf:sh-theta:P10}: \lref{GlobalTheory.T.fockVectors_pureTensors}. \ref{pf:sh-theta:P11}:
\lref{GlobalTheory.T.definiteThetaLift_localVector}. \ref{pf:sh-theta:P12}:
\lref{GlobalTheory.T.thetaChainStepVector_localPureTensor}.
\end{proof}

\subsubsection{The finite zeta integral of a theta lift}

We use the lift data of \cref{pf:onto3:adelcdef}: for $i<m$, a type $t$, an embedding $\sigma$, a local level $P$ (subgroups
$H_v$ of $\GL_r(k_v)$ with characters at the level places $v$), conjugate characters $c_\mu,c_\eta$, a Weil representation $\omega$
of $(V_c(i),W_c(i))$ ($\rank V_c(i)=n_V=r+t$, $\rank W_c(i)=r$), orthonormal bases, the \emph{level vector} $f$ (a theta function
which is the lattice vacuum outside the level and twist places and given by the local level data there) and coefficient data. Let
$S_t$ be the finite set of places of $F$ below the level places and the twist places of $t$. For an eigenvalue datum $\mathrm{ev}$ let
the \emph{type part} be the space of cusp forms in the multiplicity space of $W_c(i)$ of level $P$ with eigenvalues $\mathrm{ev}$
outside $S_t$.

\begin{theorem}[Finite zeta integral]\label{pf:sh-theta:finzeta}
Assume rapid decay of cusp forms and the unramified doubling computation for $\GL_n$. Let $R$ be a real datum and $F$ a face, and
take lift data,
a covolume-one measure on the symplectic space and an eigenvalue datum $\mathrm{ev}$ with non-zero type part. Then there is a family
$B$ of Satake parameters of rank $r$ outside $S_t$ for $\mathrm{ev}$, with $|\beta|\le\Nm w^{(r-1)/2}$ for $\beta\in B(w)$, such that
the partial $L$-series of $B$ twisted by $c_\mu^{-1}$ converges at $(1+t)/2$, and for all $U_1$ in the type part, $U_2$ in the cusp
multiplicity space and smooth vectors $v_1,v_2$,
\[
\sum_{y\in G_{\mathrm f}/K}\langle\omega(1,y)f,f\rangle\,\overline{\mathrm{pet}\bigl(y\cdot U_1(v_1),U_2(v_2)\bigr)}
=|H|\,\|f\|^2\,\frac{L^{S_t}\bigl(\frac{1+t}2,B\otimes c_\mu\bigr)}{d_{\mathrm{off}}}\,\overline{\mathrm{pet}\bigl(U_1(v_1),U_2(v_2)\bigr)},
\]
the sum over the cosets of the level subgroup $K$ of $\Uni(W_c(i))(\A_{\mathrm f})$ (an unconditionally convergent sum), where $|H|$ is the order of
the level group and $d_{\mathrm{off}}$ the normalising factor at $S_t$. That is, $\sfname{FiniteZetaL}(R,F)$ holds.
\formallabel{GlobalTheory.T.finiteZetaIntegral_thetaLift_eq_partialLValue}
\formallabel{GlobalTheory.FinZeta.finiteZetaL}
\end{theorem}

\begin{proof}
The level vector is a theta function, the vacuum off the level places
(\lref{GlobalTheory.T.levelVector_finiteDatum_isTheta_vacuum_off}), and transforms under the level subgroup by the level
character (\lref{GlobalTheory.T.finiteLevelDatumVector_transforms_by_levelCharacter}). The level of the datum is hyperspecial
outside $S_t$ and determined by its components at $S_t$ (\lref{GlobalTheory.T.datumLevel_hyperspecial_off_finite_places}), the
index of the principal level in the lift level is $|H|$ (\lref{GlobalTheory.T.principalLevel_relIndex_liftLevel_eq_card}), the
characters are unramified outside $S_t$ (\lref{GlobalTheory.T.conjugatedThetaCharacters_restriction_ramification}), and
$\langle\omega(1,g)f,f\rangle=0$ for $g$ outside the saturation of the level, because the support of the translate is moved off that
of $f$ (\lref{GlobalTheory.T.thetaTranslate_orthogonal_of_support_moved_off} and \lref{GlobalTheory.T.mem_level_of_levelDatum_ne_zero_translate}).
A non-zero element of the type part has a non-zero value at some vector; rapid decay bounds the values, and an eigenfunction with
eigenvalues $\mathrm{ev}$ yields Satake parameters $B$ with $|\beta|\le\Nm w^{(r-1)/2}$
(\lref{GlobalTheory.T.exists_satake_parameters_bounded_eigenfunction}); the parameters are then those of every element of the
type part. Since $r+5\le t$, the twisted partial $L$-series converges at $(1+t)/2$, whose real part exceeds $(r-1)/2+1$. If
$U_1(v_1)=0$ both sides vanish. Otherwise the identity is \lref{GlobalTheory.FinZeta.zeta_sum}: under these conditions (vacuum
off $S_t$, the transformation law, the vanishing off the saturation, the bounded parameters and the convergence) the coset sum
factors into local unramified zeta sums -- computed at the inert places by \cref{pf:sh-theta:inert} and at the split places by
$\sfname{UnramifiedDoublingGL}$ -- whose Euler product is the partial $L$-value divided by the normalising factor, times
$|H|\,\|f\|^2$ (about $1\,900$ Lean lines).
\end{proof}

\subsubsection{Level index and norm of the level vector}

We first record an elementary count.

\begin{lemma}[A parabolic index over a finite field]\label{pf:sh-theta:glindex}
Let $k$ be a finite field with $q$ elements. Then $|\GL_n(k)|=q^{n^2}\prod_{i=1}^n(1-q^{-i})$ for every $n\in\N$, and if
$n=r+t$,
\[
\frac{|\GL_n(k)|}{|\GL_t(k)|\,q^{rt}}\;q^{-(r^2+rt)}=\prod_{i=t+1}^{n}\bigl(1-q^{-i}\bigr).
\]
\formallabel{GlobalTheory.T.card_GL_div_parabolic_index_finiteField}
\formallabel{GlobalTheory.Twist.GLIndexLoc.index_loc}
\formallabel{GlobalTheory.Twist.GLIndexLoc.card_GL_real_eq}
\end{lemma}

\begin{proof}
An invertible $n\times n$ matrix is a sequence of $n$ columns such that the $(i+1)$-st column is not in the span of the first
$i$, which has $q^i$ elements. Hence
$|\GL_n(k)|=\prod_{i=0}^{n-1}(q^n-q^i)=q^{n^2}\prod_{i=0}^{n-1}(1-q^{i-n})=q^{n^2}\prod_{i=1}^{n}(1-q^{-i})$. Since $q\ge2$, all
factors $1-q^{-i}$ ($i\ge1$) are non-zero. Using the formula for $n$ and $t$, the left-hand side of the claim is
\[
q^{\,n^2-t^2-rt-r^2-rt}\,\frac{\prod_{i=1}^{n}(1-q^{-i})}{\prod_{i=1}^{t}(1-q^{-i})}=\prod_{i=t+1}^{n}(1-q^{-i}),
\]
because $n^2=(r+t)^2=t^2+2rt+r^2$.
\end{proof}

For a place $v$ of $F$ let $\epsilon_v=1$ if $v$ splits in $E$ and $\epsilon_v=-1$ otherwise, and
\[
d_v=\prod_{i=t+1}^{n_V}\bigl(1-\epsilon_v^i\Nm v^{-i}\bigr)^{-1}.
\]
For the lift data above let $K_{\mathrm{lev}}$ be the level subgroup of $\Uni(V_c(i))(\A_{\mathrm f})$ attached to $(t,P)$ and $K_\Lambda$
the stabiliser of the lattice of $V_c(i)$.

\begin{theorem}[Level index times norm]\label{pf:sh-theta:levelindex}
Assume the single-block local integral statement $\sfname{LocalIntegralSingleBlock}$ and $\sfname{PlaceModels}(D)$. Then for all lift
data, $K_{\mathrm{lev}}\subseteq K_\Lambda$, and for every measure of covolume one on the symplectic space
\[
[K_\Lambda:K_{\mathrm{lev}}]\cdot\|f\|^2=\prod_{v\in S_t}d_v^{-1}.
\]
That is, $\sfname{LevelIndexNormC}(R,F)$ holds.
\formallabel{GlobalTheory.T.thetaLevelIndex_mul_schwartzNorm_formula}
\formallabel{GlobalTheory.LevelIndex.levelIndexNormC}
\end{theorem}

\begin{proof}
The inclusion holds because the frame of $V_c(i)$ is a base frame of its lattice
(\lref{GlobalTheory.T.level_le_stabilizer_selfDual_lattice_frame}). Let $P_g$ be the set of level places and $Q(t)$ that of the
twist places (places of $E$), and write $k_w$ for the residue field at $w$ and $q_w=|k_w|$. We use three computations:
\begin{align}
[K_\Lambda:K_{\mathrm{lev}}]&=\prod_{w\in P_g}\frac{|\GL_{n_V}(k_w)|}{|H_w|\,|\GL_t(k_w)|\,q_w^{rt}}
\prod_{w\in Q(t)}\frac{|\GL_{n_V}(k_w)|}{|\GL_r(k_w)|\,|\GL_t(k_w)|\,q_w^{rt}},\label{eq:pf:sh-theta:index}\\
\|f\|^2&=\prod_{w\in P_g}|H_w|\,q_w^{-(r^2+rt)}\prod_{w\in Q(t)}|\GL_r(k_w)|\,q_w^{-(r^2+rt)},\label{eq:pf:sh-theta:norm}
\end{align}
(\lref{GlobalTheory.T.relIndex_level_selfDual_stabilizer_eq_prod_local}: a computation of the index place by place at the places
of $P_g\cup Q(t)$; \lref{GlobalTheory.T.levelVector_squareNorm_eq_prod_localVolumes}: from the place models and the single-block
local integral), and: $P_g$ and $Q(t)$ are disjoint, $w\mapsto w\cap\OO_F$ is
injective on $P_g\cup Q(t)$ with image $S_t$, and every $w$ in it lies over a split place $v$ with $q_w=\Nm v$
(\lref{GlobalTheory.T.levelAndTwistPlaces_split_distinct}). Multiplying \eqref{eq:pf:sh-theta:index} and
\eqref{eq:pf:sh-theta:norm}, the factors $|H_w|$ and $|\GL_r(k_w)|$ cancel, and every $w\in P_g\cup Q(t)$ contributes
\[
\frac{|\GL_{n_V}(k_w)|}{|\GL_t(k_w)|\,q_w^{rt}}\,q_w^{-(r^2+rt)}=\prod_{i=t+1}^{n_V}\bigl(1-q_w^{-i}\bigr)
\]
by \cref{pf:sh-theta:glindex}. As $v=w\cap\OO_F$ is split, $\epsilon_v=1$ and this is $d_v^{-1}$. Reindexing the product by
$S_t$ gives the claim.
\end{proof}

% ---------------------------------------------------------------------------
\subsection{Remarks on the formal statements}\label{sec:pf:sh-theta:remarks}

\begin{remark}[Unused hypotheses]\label{pf:sh-theta:rem:unused}
\begin{enumerate}[label=(\alph*)]
\item The hypothesis $F=\sfname{ThetaLocalInputs}$ of the first formal name of \cref{pf:sh-theta:weilconv} is not used; only
$\sfname{SiegelSetUnitary}$ is.
\item The hypothesis $\sfname{BernsteinZelevinskyFiniteness}$ of \cref{pf:sh-theta:finzeta} (the second of its three formal
hypotheses) is not used.
\item Of the bundle $F$, the proof of \cref{pf:sh-theta:cusp} uses only the discrete series cover, Harish-Chandra modules, rapid
decay of cusp forms and Casselman--Wallach globalisation.
\item Statement (T5) of \cref{pf:sh-theta:five} is formally stated under $F$ and Theorem~K, but these hypotheses are not used: the
statement holds for all data, with all thresholds equal to $0$.
\item In (T1), and hence in the proof of \cref{pf:sh-theta:towerthm}, the non-vanishing of the squares is not used: every square
constant equals $\kappa(t)c(t)$. Of $\sfname{LNe}$ only the consequences listed in \cref{pf:sh-theta:towerrem} are used.
\end{enumerate}
\end{remark}

\begin{remark}[Restatements]\label{pf:sh-theta:rem:wrappers}
The three formal names of \cref{pf:sh-theta:sqcor} and the second formal name of \cref{pf:sh-theta:cusp} are restatements, with the
same hypotheses and conclusions, of \cref{pf:sh-theta:sqlmthm}, \lref{GlobalTheory.SquareAgree.twisted_theta_integrals_agree},
\lref{GlobalTheory.RowClaim.rowSum_of_hypotheses_alt} and \cref{pf:sh-theta:cusp}; similarly the first formal names of
\cref{pf:sh-theta:jump,pf:sh-theta:hub,pf:sh-theta:setpt,pf:sh-theta:classsum} restate the second ones.
\end{remark}

\begin{remark}[On the form of the statements]\label{pf:sh-theta:rem:form}
\begin{enumerate}[label=(\alph*)]
\item In $\sfname{Tower}$ the numbers $u_k$ are \emph{step multiples of $1$}, i.e.\ of the form $\shthetaemb(b\prod_vg_v/g'_v)$ with
$b$ a $Q(t)$-unit and $g_v,g'_v$ divisors of $\Nm v^{c(v)}$, $v\in S_0$; they need not be units at the places of $S_0$. An
element $x$ of $\shthetaLL$ is an $n$-unit if $xn^k$ is an integer dividing a power of $n$; the definition does not require
$x\ne0$ (for $n\ge1$ it excludes $x=0$).
\item $\sfname{LiftsCuspidal}$ is a statement about data with $r>1800$ and Harish-Chandra parameters at distance more than
$2N(D)^3$; this size condition is part of the statement.
\item In the application of \cref{pf:sh-theta:hub} in the proof of \cref{pf:sh-theta:setpt}, the constant $C_T$ is $0$, so that the
inequality of (H7) only says $0<g_k$, and only the rest point conditions of (H1) are used.
\item \Cref{pf:sh-theta:inert,pf:sh-theta:placemodels} have no hypotheses from the literature.
\item The formal name of \cref{pf:sh-theta:setpt} refers to places which are ``not twisted'' (\cref{pf:sh-theta:setptrem}).
\end{enumerate}
\end{remark}
